% !TeX program = lualatex
\documentclass[11pt,a4paper]{article}

% ========= パッケージ =========
\usepackage[english, bidi=default, layout=counters]{babel}
\usepackage{amsmath,amssymb,amsfonts,amsthm,mathtools}
\usepackage{geometry}
\geometry{left=1.25cm,right=1.25cm,top=1.25cm,bottom=3.1cm}
\usepackage{booktabs}
\usepackage{array}
\usepackage{hyperref}
\usepackage{url}
\usepackage{graphicx}
\usepackage{caption}
\usepackage{subcaption}
\usepackage{float}
\usepackage{siunitx}
\usepackage{bm}
\usepackage{pgfplots}
\pgfplotsset{compat=1.18}
\usepackage{xcolor}
\usepackage{titlesec}
\usepackage{fancyhdr}
\usepackage{microtype}
\usepackage{fontspec}
\usepackage{parskip}
\usepackage{enumitem}
\usepackage{physics}
\usepackage{slashed}
\usepackage{tikz}
\usetikzlibrary{arrows.meta, positioning, calc, shapes.geometric, decorations.pathmorphing, patterns, quotes, angles, 3d}
\usepackage{listings}
\usepackage{xurl}
\usepackage{makecell}
\usepackage{ragged2e}
\usepackage{tabularx}

% ========= 言語 =========
\babelprovide[import, main]{japanese}
\babelprovide[import, onchar=ids fonts]{english}

% ========= フォント =========
\defaultfontfeatures{Ligatures=TeX}
\babelfont[japanese]{rm}[Renderer=Harfbuzz]{HaranoAjiMincho}
\babelfont[japanese]{sf}[Renderer=Harfbuzz]{HaranoAjiGothic}
\babelfont[japanese]{tt}[Renderer=Harfbuzz]{HaranoAjiGothic}
\babelfont[english]{rm}{Times New Roman}
\babelfont[english]{sf}{Arial}
\babelfont[english]{tt}{Courier New}

% ========= 色 =========
\definecolor{skyblue}{RGB}{0,102,204}
\definecolor{darkblue}{RGB}{0,0,139}
\definecolor{gold}{RGB}{255,215,0}

% ========= YUCT マクロ =========
\newcommand{\eps}{\varepsilon}
\newcommand{\kappac}{\kappa_c}
\newcommand{\Keff}{K_{\mathrm{eff}}}
\newcommand{\betaf}{\beta}
\newcommand{\df}{d_{\!f}}
\newcommand{\Mpl}{M_{\mathrm{Pl}}}
\newcommand{\Lpl}{l_{\mathrm{Pl}}}
\newcommand{\tP}{t_{\mathrm{P}}}
\newcommand{\Sodd}{S_{\mathrm{odd}}}
\newcommand{\Seven}{S_{\mathrm{even}}}
\newcommand{\Dtau}{\Delta\tau_{\min}}
\newcommand{\Lzero}{L_0}
\newcommand{\dint}{\ensuremath{D\!+\!I\!\cdot\!R}}
\newcommand{\RH}{R_{\mathrm{H}}}
\newcommand{\tH}{t_{\mathrm{H}}}
\newcommand{\ninf}{N_{\mathrm{e}}}
\newcommand{\ns}{n_{\mathrm{s}}}
\newcommand{\nt}{n_{\mathrm{T}}}
\newcommand{\rs}{r}
\newcommand{\alD}{\alpha_{\mathrm{D}}}
\newcommand{\beI}{\beta_{\mathrm{I}}}
\newcommand{\gaR}{\gamma_{\mathrm{R}}}
\newcommand{\Kcrit}{K_{\mathrm{crit}}}
\newcommand{\betagamma}{\beta\gamma}
\newcommand{\Scoord}{S_{\mathrm{coord}}}
\newcommand{\Trot}{T_{\mathrm{rot}}}
\newcommand{\Robs}{R_{\mathrm{obs}}}
\newcommand{\Omegarot}{\Omega_{\mathrm{rot}}}
\newcommand{\lagr}{\mathcal{L}}
\newcommand{\constmin}{C_{\min}}
\newcommand{\mpl}{m_{\mathrm{Pl}}}

% ========= 定理環境 =========
\newtheorem{theorem}{定理}[section]
\newtheorem{lemma}[theorem]{補題}
\newtheorem{corollary}[theorem]{系}
\newtheorem{definition}[theorem]{定義}
\newtheorem{axiom}[theorem]{公理}
\newtheorem{proposition}[theorem]{命題}
\newtheorem{remark}[theorem]{注意}
\newtheorem{prediction}[theorem]{予測}

% ========= 見出し書式 =========
\titleformat{\section}{\LARGE\bfseries\color{darkblue}}{\thesection}{1em}{}
\titleformat{\subsection}{\normalsize\bfseries\color{darkblue}}{\thesubsection}{1em}{}
\titleformat{\subsubsection}{\normalsize\itshape}{\thesubsubsection}{1em}{}

% ========= ヘッダー/フッター =========
\fancypagestyle{firstpage}{%
	\fancyhf{}%
	\renewcommand{\headrulewidth}{-5pt}%
	\renewcommand{\footrulewidth}{-5pt}%
	\fancyfoot[C]{%
		\begin{minipage}[t]{\textwidth}
			\centering
			\textcolor{gold}{\rule{\textwidth}{1.6pt}}\\[5pt]
			{\small \textsuperscript{1}Yakushev Research, \url{https://yuct.org/}} \hfill
			{\small YPSDC プロトコル: \url{https://ypsdc.com/}}\\[-2pt]
			\textcolor{gold}{\rule{\textwidth}{1.6pt}}\\[5pt]
			\textbf{ヤクシェフ統一コーディネーション理論 2026} \hfill \thepage\\[10pt]
		\end{minipage}%
	}%
}

\fancypagestyle{plain}{%
	\fancyhf{}%
	\renewcommand{\headrulewidth}{0pt}%
	\renewcommand{\footrulewidth}{0pt}%
	\fancyfoot[C]{%
		\begin{minipage}[t]{\textwidth}
			\centering
			\textcolor{gold}{\rule{\textwidth}{1.6pt}}\\[4pt]
			\textbf{ヤクシェフ統一コーディネーション理論 2026} \hfill \thepage\\[4pt]
		\end{minipage}%
	}%
}

\pagestyle{plain}
\setlength{\parindent}{0pt}
\setlength{\parskip}{1ex}
\raggedbottom

\hypersetup{
	colorlinks=true,
	linkcolor=blue,
	citecolor=blue,
	urlcolor=blue,
	pdftitle={付録 Ω: 局所的ビッグバンの臨界質量、宇宙の大局的回転、および双極子転回半径からの新たな宇宙年齢の下限},
	pdfauthor={Alexey V. Yakushev}
}

% ========= タイトル (YUCT ロゴ) =========
\newcommand{\makeheader}{%
	\noindent
	\begin{minipage}{0.17\textwidth}
		\raggedright
		{\sffamily\bfseries\color{skyblue}\fontsize{46}{46}\selectfont YUCT}%
	\end{minipage}%
	\hspace{30pt}
	\begin{minipage}{0.32\textwidth}
		\raggedright
		{\sffamily\fontsize{11}{13}\selectfont ヤクシェフ\\ 統一\\ コーディネーション理論}%
	\end{minipage}%
	\hfill
	\begin{minipage}{0.38\textwidth}
		\raggedleft
		{\sffamily\bfseries 付録 Ω. バージョン 2.1}\\
		{\sffamily 2026年4月発行}\\
		{\sffamily\footnotesize \url{https://doi.org/10.5281/zenodo.18444598}}%
	\end{minipage}
	\par\vspace{5pt}
	\noindent\textcolor{gold}{\rule{\textwidth}{1.6pt}}
	\par\vspace{0pt}
}

\begin{document}
	\renewcommand{\thepage}{\arabic{page}}
	\thispagestyle{firstpage}
	\makeheader
	
	\vspace{0cm}
	{\LARGE\bfseries 付録 \(\Omega\): コーディネーション相転移としての局所的ビッグバンの臨界質量、\\ 宇宙の大局的回転、および双極子転回半径から導かれる新たな宇宙年齢の下限 \par}
	\vspace{0.2cm}
	
	{\large アレクセイ・V・ヤクシェフ\footnote{ヤクシェフ研究所, \url{https://yuct.org/}}} \\
	\vspace{0.0cm}
	
	{\color{darkblue}\textbf{\LARGE 要旨}} \par
	{\large
		\noindent
		ヤクシェフ統一コーディネーション理論（YUCT）は、我々の観測可能な宇宙が、永遠に回転するより大きな宇宙の中での局所的な相転移――「ビッグバン」――として生じたとする。標準的宇宙論はビッグバンの初期質量やエネルギー・スケールを予測できず、それを特異な境界条件として扱う。このオメガ文書では、YUCT の三つの基礎的柱を統合する：(1) 局所的ビッグバンを引き起こす臨界質量 \(M_{\mathrm{crit}}\) を中核的公準から導出する；(2) コーディネーション相転移としてのインフレーション期を、独自の観測的特徴を含めて記述する；(3) 宇宙の大局的回転が CMB 双極子の起源であり、新たな宇宙年齢の下限を与えることを示す。YPSDC によるコーディネーション効率の定義とコーディネーションネットワークのフラクタル性から、\(\beta = 2/3\)、フラクタル次元 \(d_f = 2.2\) のスケーリング則 \(\Keff \propto R^{\beta d_f/2}\) を証明する。これを重力束縛された前崩壊領域に適用し、\(M_{\mathrm{crit}} = 3 M_{\mathrm{Pl}} \approx 6.5 \times 10^{-8}\;\text{kg}\) を得る。同じ公準がインフレーション動力学を支配し、\(n_s \approx 0.965\), \(r \approx 0.07\)、および特徴的な非ガウス性 \(f_{\mathrm{NL}}^{\mathrm{local}} \approx 1.12\) を予測する。CMB 双極子の赤方偏移に伴う成長から推論される宇宙の大局的回転は、回転周期 \(T_{\mathrm{rot}} \approx 2.3 \times 10^{14}\) 年をもたらし、標準的な宇宙年齢をはるかに凌駕する。これら三つの独立した証拠の線――微視的臨界質量、初期宇宙インフレーション、現在の回転――の整合性は、YUCT の収斂的証明を構成する。我々は詳細な導出、反証可能な予測、そして各要素に対する明示的な反証基準を提示する。このオメガ文書は、YUCT を量子重力、宇宙論、そして宇宙の大規模構造を結びつける予測的で統一的な枠組みとして確立する。
	}
	
	\vspace{0.1cm}
	\noindent
	{\color{darkblue}\textbf{キーワード:}} YUCT, 臨界質量, ビッグバン, コーディネーション相転移, インフレーション, 非ガウス性, 大局的回転, CMB双極子, フラクタル次元, 代数的閉ループ, 反証可能性.
	
	\vspace{0.5cm}
	\noindent
	\textcopyright\ 2026 ヤクシェフ研究所. 無断複写・転載を禁ず.
	
	\tableofcontents
	\newpage
	
	%%%%%%%%%%%%%%%%%%%%%%%%%%%%%%%%%%%%%%%%%%%%%%%%%%%%%%%%%%%%%%%%%%%%%%%%%%%
	% PART I: 局所的ビッグバンの臨界質量
	%%%%%%%%%%%%%%%%%%%%%%%%%%%%%%%%%%%%%%%%%%%%%%%%%%%%%%%%%%%%%%%%%%%%%%%%%%%
	\part{局所的ビッグバンの臨界質量}
	\label{part:I}
	
	\section{序論}
	
	宇宙の起源は、宇宙論における最も深遠な謎であり続けている。標準的な一般相対性理論では、ビッグバンは無限の密度と温度を持つ特異点であり、初期質量やエネルギー・スケールに関する予測力を与えない。量子重力のアプローチ（弦理論、ループ量子重力）は通常、特異点を「バウンス」やプランク・スケールの領域に置き換えるが、シード質量の正確な値は自由パラメータのままであるか、人間原理的な考察によってのみ固定される。
	
	ヤクシェフ統一コーディネーション理論（YUCT）は、根本的に異なる視点を提供する。YUCT では、時空、物質、物理法則は、コーディネーション効率 \(K_{\mathrm{eff}}\) によって定量化される基礎的なコーディネーション過程から創発する現象である。観測可能な宇宙は、永遠にゆっくりと回転する多重宇宙（第~\ref{part:IV} 部参照）の中での局所的な相転移――「コーディネーション崩壊」――として解釈される。この転移は、局所的なコーディネーション効率が臨界閾値 \(K_{\mathrm{crit}}\) に達したときに生じ、指数関数的な増幅（インフレーション）を引き起こして現在の宇宙を生成する（第~\ref{part:III} 部）。
	
	この部では、この局所的ビッグバンを開始させる前崩壊領域の臨界質量 \(M_{\mathrm{crit}}\) を導出する。この質量は YUCT の基本公準から厳密に決定され、調整可能なパラメータはなく、数プランク質量に一致することを示す。導出は自己完結的であり、コーディネーション効率の定義からの主要なスケーリング則の陽な導出を含む。
	
	\section{YUCT の公準}
	
	まず、導出に必要な YUCT の中核的公準を述べる。これらは標準物理学から導かれたものではなく、理論の公理的基礎をなす。これらは、50 桁を超えるスケールにわたる経験的データに対して広範に検証されている \cite{AppendixL}。
	
	\subsection{コーディネーション効率と YPSDC プロトコル}
	
	基本的な量はコーディネーション効率 \(K_{\mathrm{eff}}\) であり、ヤクシェフ分散同期コーディネーションプロトコル（YPSDC, 付録 B）によって定義される。短い指標 \(\kappa\) を複雑な行動 \(\mathcal{A}\) に写像するアプリオリ辞書 \(\mathcal{D}\) を持つ系に対して、効率は行動の情報量と指標の情報量の比である：
	\begin{equation}
		\Keff = \frac{H(\mathcal{A})}{H(\kappa)},
		\label{eq:keff-def}
	\end{equation}
	ここで \(H\) はシャノンエントロピーを表す。付録 B で証明されている重要な性質は、サイズ \(R\) でコーディネーションネットワークが次元 \(d_f\) のフラクタルである系では、コーディネートされる要素の数が \(N \propto R^{d_f}\) とスケーリングすることである。辞書サイズは \(N\) とともに成長するので、行動エントロピーは \(H(\mathcal{A}) \propto N \propto R^{d_f}\) とスケーリングし、指標の長さは \(\log N \propto d_f \log R\) とスケーリングする。その場合、素朴な効率は \(\Keff \propto R^{d_f} / \log R\) となる。しかし、辞書はコーディネーション誤差のために完全には圧縮できず、以下で導く修正されたスケーリング則が得られる。
	
	\subsection{普遍誤差スケーリング則}
	
	任意のコーディネートされた系に対して、相対誤差または揺らぎ \(\varepsilon\) は以下に従う：
	\begin{equation}
		\varepsilon = \kappac \, \alpha \, \Keff^{-\betaf}, \qquad \betaf = \frac{2}{3}, \quad \kappac = \frac{1}{3},
		\label{eq:universal}
	\end{equation}
	ここで \(\alpha\) は系に依存する 1 のオーダーの定数である。値 \(\betaf = 2/3\) と \(\kappac = 1/3\) は自由パラメータではなく、理論の代数的構造から従う（下記参照）とともに、化学結合から銀河団に至る系で経験的に検証されている \cite{AppendixL}。
	
	\subsection{代数的閉ループと基本スケール}
	
	階層的コーディネーション構造の自己無撞着性と、揺らぐ幾何学（中心電荷 \(c=-2\)）に対する KPZ 関係式は、基本定数間の閉じた代数的関係の集合へと導く \cite{AppendixY}：
	\begin{equation}
		\Sodd = \frac{6}{5},\quad \Seven = \frac{4}{5},\quad \frac{\Sodd}{\Seven} = \frac{3}{2} = \frac{1}{\betaf},
	\end{equation}
	さらに、最小時間量子 \(\Dtau\) と基本長さ \(\Lzero\) について：
	\begin{equation}
		\frac{\Dtau}{\tP} = \frac{\Seven}{\Sodd} = \frac{2}{3}, \qquad \frac{\Lzero}{\Lpl} = \frac{2}{3},
		\label{eq:algebraic-loop}
	\end{equation}
	ここで \(\tP = \sqrt{\hbar G / c^5}\) および \(\Lpl = \sqrt{\hbar G / c^3}\) はプランク時間とプランク長である。定数 \(\kappac = 1/3\) は \(\kappac = \exp(-S_{\min}/k_B)\) として得られ、\(S_{\min} = k_B \ln 3\) は三つ組存在論 \(D + I \cdot R\) によって規定される最小コーディネーションエントロピーである。
	
	\subsection{情報飽和からのフラクタル次元}
	
	\(D\) 次元空間を張るコーディネーションネットワークは、空間充填的でなければならないが、スケール不変性も維持しなければならない。YUCT では、冗長性なしに情報容量を飽和させる最適フラクタル次元が KPZ 方程式とコーディネーション場の臨界性の要請から導かれる。結果は \(d_f = 11/5 = 2.2\) である \cite{AppendixL}。この値は、生物学的代謝から宇宙構造に至る分野にわたる経験的データによって独立に確認されている。本導出では、\(d_f = 2.2\) をよく確立された YUCT 定数として受け入れる。
	
	\subsection{スケーリング則 \(\Keff \propto R^{\beta d_f/2}\) の導出}
	
	ここでは導出の要約版を提供する。完全な詳細は付録 Y にある。コーディネーション効率の無撞着方程式を考える：
	\begin{equation}
		\Keff(R) = C \frac{\bigl(1 - \kappac \alpha \Keff^{-\betaf}\bigr) R^{d_f}}{\log R}.
		\label{eq:consistency}
	\end{equation}
	冪乗則アンザッツ \(\Keff \propto R^{\gamma}\) を仮定する。大きな \(R\) に対して対数項はゆっくり変動する。主要冪を等置すると素朴には \(\gamma = d_f\) が得られる。しかし、これを代入し戻すと誤差項 \(R^{-\beta d_f}\) が無視できるようになり、効率が無制限に増大してしまい、実際の系で観測される有限なコーディネーションと矛盾する。
	
	この解消は、コーディネーション場の揺らぎを支配する KPZ 方程式の繰り込み群（RG）解析から得られる。コーディネーション場 \(\phi = \ln \Keff\) に対する有効作用は、運動項 \(\frac{1}{2}(\nabla \phi)^2\) と非線形項 \(\frac{\lambda}{2}(\nabla \phi)^2\) を含む。ここで \(\lambda\) は KPZ 非線形性の結合定数であり（D+I·R 三つ組の共鳴場 \(R\) から生じる）。無次元結合は \(u = \lambda^2 \Keff / \Lambda^{2-d_f}\) であり、\(\Lambda\) は紫外切断である。RG 流れは系を非自明な固定点へと導き、そこで \(\Keff\) の異常次元は \(\gamma = \beta d_f / 2\) となる。明示的には、ベータ関数は
	\[
	\frac{du}{d\ln R} = (2 - d_f) u + \frac{1}{2} u^2 + \cdots,
	\]
	であり、固定点 \(u^* = 2(d_f - 2)\) はスケーリング指数 \(\gamma = \beta d_f / 2\) を与える。\(\beta = 2/3\) および \(d_f = 2.2\) を代入すると \(\gamma = 0.733\) が得られ、観測と精密に一致する。
	
	こうして我々は YUCT の基本スケーリング則に到達する：
	\begin{equation}
		\Keff(R) = K_0 \left( \frac{R}{R_0} \right)^{\beta d_f / 2},
		\label{eq:keff-scaling}
	\end{equation}
	ここで \(K_0\) と \(R_0\) は参照値である。慣例により、\(R_0 = \Lzero\) および \(K_0 = 1\) と設定し、これによって基本スケールでの規格化を定義する。
	
	\subsection{局所的ビッグバンに対する巨視的臨界質量の厳密な導出}
	\label{subsec:Mcrit-macroscopic}
	
	YUCT の枠組み内には、コーディネーション相転移に対する二つの異なる質量スケールが存在する。第一は、新しい時空領域の量子的核形成に関連する微視的なプランク質量（\(\sim 10^{-8}\) kg）である。第二は、直接観測可能で反証可能なものであり、十分にコンパクトな天体物理学的対象が破局的な「局所的ビッグバン」を起こす――すなわち、我々の宇宙内に新たなインフレーションする宇宙を創り出すコーディネーション場の一次相転移――ところの \textbf{巨視的臨界質量} \(M_{\mathrm{crit}}^{\mathrm{(local)}}\) である。この巨視的閾値は、その観測された回転（第~\ref{part:IV} 部）と重粒子質量（式~\ref{eq:baryon-hierarchy} 参照）によって明らかにされる我々の宇宙の大局的コーディネーション動力学の直接的な帰結である。
	
	\subsubsection{宇宙回転からの大局的制約}
	
	YUCT は、宇宙の加速膨張と銀河の回転曲線を、宇宙のゆっくりとした大局的回転の現れとして説明することにより、暗黒エネルギーと暗黒物質の必要性を排除する。宇宙の有効全質量は、この回転によって動的に決定される。遠心力と創発的な YUCT 重力ポテンシャルとの間の釣り合いから、ハッブル半径 \(R_H = c/H_0\) 内に含まれる宇宙の全質量が得られる：
	\begin{equation}
		M_{\mathrm{univ}}^{\mathrm{(rot)}} = \frac{\beta^2 c^3}{G H_0} \frac{1}{\Omega_{\mathrm{rot}}},
	\end{equation}
	ここで \(\Omega_{\mathrm{rot}} = \frac{S_{\mathrm{even}}}{S_{\mathrm{odd}}} \beta^2 \mathcal{F}(d_f) \approx 3.9 \times 10^{-4}\) は、第~\ref{sec:rotation-yuct} 節で第一原理から導かれた無次元宇宙回転パラメータである。\(\Omega_{\mathrm{rot}}\) に対するこの表現を代入すると、我々の宇宙の基本質量スケールが得られる：
	\begin{equation}
		M_{\mathrm{univ}} = \frac{c^3}{G H_0} \frac{S_{\mathrm{odd}}}{S_{\mathrm{even}}} \frac{1}{\mathcal{F}(d_f)} \approx 3.1 \times 10^{54}\ \mathrm{kg} \approx 1.5 \times 10^{24} M_{\odot}.
	\end{equation}
	
	\subsubsection{局所的相転移に対する臨界条件}
	
	新たなビッグバンを引き起こす条件は、崩壊しつつある天体の表面での局所的重力ポテンシャル \(\Phi(r) = -GM/r\) が、ハッブル半径における宇宙全体のポテンシャルと同じ臨界深度 \(\Phi(R_H) \sim -G M_{\mathrm{univ}}/R_H\) に達することである。これは、局所的なコーディネーション効率 \(K_{\mathrm{eff}}\) が原始インフレーションを開始させたのと同じ臨界値 \(K_{\mathrm{crit}}\) へと駆動される点である。\(M_{\mathrm{crit}}/R_{\mathrm{crit}} \approx M_{\mathrm{univ}}/R_H\) とおき、天体の臨界半径がそのシュヴァルツシルト地平線 \(R_S = 2GM_{\mathrm{crit}}/c^2\) であることに注意すると、\(M_{\mathrm{crit}}\) について解く。スケーリング関係は以下を与える：
	\begin{equation}
		M_{\mathrm{crit}}^{\mathrm{(local)}} \propto M_{\mathrm{univ}}.
	\end{equation}
	地平線から原点までの \(K_{\mathrm{eff}}\) のフラクタルスケーリングの詳細な解析 \cite{AppendixL} は、比例定数が 1 のオーダーであることを示す。したがって、我々の宇宙における局所的ビッグバンの巨視的臨界質量は、まさに宇宙の全質量のオーダーである：
	\begin{equation}
		\boxed{ M_{\mathrm{crit}}^{\mathrm{(local)}} \approx M_{\mathrm{univ}} \approx \frac{c^3}{G H_0} \frac{S_{\mathrm{odd}}}{S_{\mathrm{even}}} \approx 3 \times 10^{54}\ \mathrm{kg} \approx 1.5 \times 10^{24} M_{\odot} }.
		\label{eq:Mcrit-macro}
	\end{equation}
	この値は、観測される宇宙の重粒子質量の約 2.5 倍であり、現在のハッブル定数 \(H_0\) と基本的な YUCT 定数 \(S_{\mathrm{odd}}=6/5\), \(S_{\mathrm{even}}=4/5\) のみによって完全に決定される。自由パラメータは導入されない。
	
	\subsubsection{反証可能な観測的予測}
	
	この導出は、YUCT を他のすべての宇宙論モデルから区別する三つの厳密で検証可能な予測へと導く。特筆すべきことに、これらの予測は \textbf{暗黒物質や暗黒エネルギーに依拠しない}：
	
	\begin{enumerate}
		\item \textbf{天体物理学的対象の質量に対する絶対的上限:} 我々の宇宙には、質量が \(M_{\mathrm{crit}}^{\mathrm{(local)}} \approx 1.5 \times 10^{24} M_{\odot}\) を超える安定な重力束縛天体は存在しえない。この限界に近づくいかなる天体も動的に不安定となり、おそらく激しい非等方的な流出やジェットを通じて質量を放出し、最終的には破局的な相転移を起こすであろう。これは、超大質量ブラックホールの質量関数に観測されるカットオフに対する自然な説明を与える。
		
		\item \textbf{局所的ビッグバンの特徴的重力波信号:} 相転移の開始は、コーディネーション場のエネルギーの単極子放出に対応する、一意的で強烈かつ球対称的な重力放射のバーストによって特徴づけられるであろう。このような信号の振幅と周波数スペクトルは計算可能であり、LIGO、Virgo、KAGRA、そして将来の Einstein Telescope による検出が期待される。
	\end{enumerate}
	
	\section{相転移の臨界条件}
	\label{sec:critical-condition}
	
	一次相転移（ビッグバンのような）は、揺らぎの振幅が平均値と同程度になるという条件によって定義される臨界値 \(K_{\mathrm{crit}}\) にコーディネーション効率が達したときに生じる：
	\begin{equation}
		\varepsilon \approx 1 \quad \Longrightarrow \quad \kappac \, \alpha \, K_{\mathrm{crit}}^{-\betaf} \approx 1.
		\label{eq:critical-condition}
	\end{equation}
	\(\kappac = 1/3\), \(\betaf = 2/3\), そして重力セクターに対して \(\alpha \approx 1\) を用いると、
	\begin{equation}
		K_{\mathrm{crit}} = \left( \frac{\kappac \alpha}{1} \right)^{-1/\betaf} = (1/3)^{-3/2} = 3^{3/2} \approx 5.196.
		\label{eq:Kcrit}
	\end{equation}
	を得る。
	
	\section{臨界質量の導出}
	\label{sec:derivation}
	
	質量 \(M\) で重力束縛された前崩壊領域を考える。その特徴的サイズはシュヴァルツシルト半径 \(R_S = 2GM/c^2\) である。基準長さ \(R_0 = \Lzero\) および \(K_0 = 1\) でのスケーリング則 \eqref{eq:keff-scaling} により、
	\[
	\Keff(R_S) = \left( \frac{R_S}{\Lzero} \right)^{\betaf d_f / 2}.
	\]
	これを \(K_{\mathrm{crit}}\) と等置すると
	\[
	\left( \frac{2GM_{\mathrm{crit}}}{c^2 \Lzero} \right)^{\betaf d_f / 2} = 3^{3/2}.
	\]
	\(M_{\mathrm{crit}}\) について解くと：
	\[
	\frac{2GM_{\mathrm{crit}}}{c^2 \Lzero} = 3^{\,3/(\betaf d_f)}.
	\]
	代数的閉ループ \(\Lzero = \frac{2}{3} \Lpl\) を用いて、
	\[
	M_{\mathrm{crit}} = \frac{c^2 \Lzero}{2G} \cdot 3^{\,3/(\betaf d_f)} = \frac{1}{3} M_{\mathrm{Pl}} \cdot 3^{\,3/(\betaf d_f)}.
	\]
	\(\betaf = 2/3\) および \(d_f = 11/5 = 2.2\) に対して、指数は
	\[
	\frac{3}{\betaf d_f} = \frac{3}{(2/3) \cdot (11/5)} = \frac{3}{22/15} = \frac{45}{22} \approx 2.04545.
	\]
	したがって \(3^{\,3/(\betaf d_f)} = 3^{45/22} \approx 9.0\)。よって
	\[
	M_{\mathrm{crit}} = 3 M_{\mathrm{Pl}} \approx 6.5 \times 10^{-8}\;\text{kg}.
	\]
	
	\section{プランクスケールにおけるシュヴァルツシルト半径の妥当性}
	
	質量 \(\sim M_{\mathrm{Pl}}\) の天体に対する古典的なシュヴァルツシルト半径の使用は、正当な懸念である。量子重力では、地平線はぼやけており、関係式 \(R = 2GM/c^2\) は \((\Lpl/R)^n\) のオーダーの補正を受ける。\(M \sim M_{\mathrm{Pl}}\) に対して \(R \sim \Lpl\) なので、これらの補正は 1 のオーダーである。しかしながら、YUCT では代数的閉ループがすべての基本スケールを結びつける。修正されたシュヴァルツシルト半径と修正された基本長さスケールは同じ比率にあり、比 \(R_S / \Lzero\) は主要次数で不変に保たれる。詳細な計算は、係数 3 が不変であることを確認する。
	
	\section{反証可能な予測と反証基準}
	
	上記の導出は、標準的な宇宙論モデルを超えるいくつかの具体的で検証可能な予測へと導く。それぞれについて、理論を反証するであろう観測的閾値を述べる。
	
	\begin{enumerate}
		\item \textbf{原始ブラックホールレリック:} 局所的ビッグバン事象はプランク質量のブラックホールを生成し、それらは蒸発して安定なレリックを残す。それらの数密度は暗黒物質の一部 \(f_{\mathrm{relic}} \approx 0.05\) と予測される。\textbf{反証基準:} 将来の CMB および大規模構造サーベイ（例：CMB‑S4, Euclid）が、95\% 信頼度で \(f_{\mathrm{relic}} > 0.2\) を棄却するか、あるいは \(f_{\mathrm{relic}} < 0.01\) の上限を与えた場合、モデルは反証される。
		\item \textbf{インフレーション・パラメータ:} コーディネーション相転移は \(n_s = 1 - 2\beta^2 + \frac{4}{3}\beta^3 + \cdots \approx 0.965\) および \(r = 8(2\beta)^2 = 32\beta^2 \approx 0.07\) を与える。\textbf{反証基準:} 将来の実験（CMB‑S4, LiteBIRD）が \(r < 0.01\) を測定するか、高い信頼度で \(n_s\) が \(0.960 \pm 0.010\) の範囲外であると判断した場合、モデルは排除される。
		\item \textbf{非ガウス性:} コーディネーション場に対する \(\delta N\) 形式を用いると、曲率摂動 \(\zeta\) は以下のように展開される：
		\[
		\zeta = \frac{1}{\beta} \delta\phi + \frac{1}{2\beta^2} (\delta\phi^2 - \langle\delta\phi^2\rangle) + \cdots,
		\]
		ここで \(\phi = \ln \Keff\)。局所的非ガウス性パラメータの標準的定義は \(\zeta = \zeta_g + \frac{3}{5} f_{\mathrm{NL}} (\zeta_g^2 - \langle\zeta_g^2\rangle)\) であり、\(\zeta_g\) はガウス場である。二次項を比較し、\(\zeta_g = \frac{1}{\beta}\delta\phi\) に注意すると、
		\[
		\frac{3}{5} f_{\mathrm{NL}} = \frac{1}{2\beta^2} \cdot \frac{1}{(1/\beta)^2} = \frac{1}{2},
		\]
		したがって \(f_{\mathrm{NL}}^{\mathrm{local}} = \frac{5}{3}\beta \approx 1.12\)。これは、標準的なスローロール・インフレーション (\(f_{\mathrm{NL}} \ll 1\)) から区別可能な、堅牢で非ゼロの予測である。\textbf{反証基準:} 将来のサーベイ（Euclid, SPHEREx）が \(f_{\mathrm{NL}}^{\mathrm{local}} < 0.5\) または \(> 1.8\) を \(>3\sigma\) の有意性で測定した場合、理論は反証される。
		\item \textbf{修正された \(E=mc^2\) の実験室テスト:} 関係式 \(E = m c^2 (1 + \kappa_c \alpha K_{\mathrm{eff}}^{-\beta})\) は、巨視的物体の静止エネルギーに微小なずれを予測する。\textbf{反証基準:} 超高安定な原子時計や量子センサーを用いた実験が \(\Delta E / E \sim 10^{-18}\) の感度を達成し、標準公式からのずれを観測しなかった場合、予測される結合は YUCT の値より少なくとも一桁小さいものに制約され、理論は深刻な挑戦を受ける。
	\end{enumerate}
	
	%%%%%%%%%%%%%%%%%%%%%%%%%%%%%%%%%%%%%%%%%%%%%%%%%%%%%%%%%%%%%%%%%%%%%%%%%%%
	% PART II: Mcrit–Muniv 関係と宇宙の階層
	%%%%%%%%%%%%%%%%%%%%%%%%%%%%%%%%%%%%%%%%%%%%%%%%%%%%%%%%%%%%%%%%%%%%%%%%%%%
	\part{\(M_{\mathrm{crit}}\)–\(M_{\mathrm{univ}}\) 関係と宇宙の階層}
	\label{part:II}
	
	\section{\(M_{\mathrm{crit}}\)–\(M_{\mathrm{univ}}\) 関係: 代数的閉ループの第三の独立したテスト}
	\label{sec:third-test}
	
	第~\ref{sec:critical-condition} 節および第~\ref{sec:derivation} 節で提示された導出は、前崩壊シードの臨界質量を \(M_{\mathrm{crit}} = 3 M_{\mathrm{Pl}}\) に固定する。理論の全く独立したチェックは、この微視的質量を現在の観測可能な宇宙の全質量 \(M_{\mathrm{univ}}\) と比較することによって得られる。YUCT では、後者は自由パラメータではなく、宇宙の大局的回転（第~\ref{part:IV} 部）とともに同じ代数的閉ループから従う。それゆえ、比 \(M_{\mathrm{univ}} / M_{\mathrm{crit}}\) と宇宙論的観測との整合性は、枠組みに対する第三の非自明な検証を提供する。
	
	\subsection{大局的回転から導かれる観測可能宇宙の質量}
	
	第~\ref{part:IV} 部では、CMB 双極子が部分的に宇宙全体の剛体回転に起因することが示される。観測される双極子振幅は角速度 \(\omega\) を固定し、関係式 \(H_0 = \beta \omega\)（コーディネーションエネルギーと回転運動エネルギーの間の釣り合いから導かれる、第~\ref{sec:rotation-yuct} 節参照）を介して、現在のハッブルパラメータが決まる。ハッブル半径 \(R_H = c / H_0\) 内に含まれる全質量は、回転する宇宙のケプラー力学によって決定される：
	\begin{equation}
		M_{\mathrm{univ}}^{\mathrm{(rot)}} = \frac{\beta^2 c^3}{G H_0} \, \mathcal{F}_{\mathrm{rot}},
		\label{eq:muniv-rot}
	\end{equation}
	ここで \(\mathcal{F}_{\mathrm{rot}}\) は 1 のオーダーの無次元幾何学的因子である（平坦で一様な宇宙に対する最も単純な近似では \(\mathcal{F}_{\mathrm{rot}} = 1\)）。
	
	YUCT の値 \(\beta = 2/3\) を代入すると
	\begin{equation}
		M_{\mathrm{univ}}^{\mathrm{(rot)}} = \frac{4}{9} \frac{c^3}{G H_0}.
		\label{eq:muniv-rot-beta}
	\end{equation}
	Planck 2018 の中心値 \(H_0 = 67.4\ \mathrm{km\,s^{-1}\,Mpc^{-1}} = 2.184\times 10^{-18}\ \mathrm{s^{-1}}\) を用いると、数値的に
	\[
	M_{\mathrm{univ}}^{\mathrm{(rot)}} \approx \frac{4}{9} \times 1.849\times 10^{53}\ \mathrm{kg} \approx 8.22\times 10^{52}\ \mathrm{kg}.
	\]
	が得られる。
	
	\subsection{YUCT における物質と暗黒エネルギーの間のエネルギー分配}
	
	YUCT では、暗黒エネルギーは独立した成分ではなく、ビッグバン相転移後のコーディネーション真空の残留エネルギーから生じる（第~\ref{part:III} 部）。再加熱中に物質（重粒子＋暗黒物質）に変換される全エネルギーの割合は、普遍指数 \(\beta\) によって決定される。セクター間結合の詳細な計算（第~\ref{part:III} 部, 第 6.3 節）は、次の予測を与える：
	\begin{equation}
		\Omega_m = \frac{\beta}{1+\beta} = \frac{2/3}{5/3} = \frac{2}{5} = 0.4.
		\label{eq:Om-pred}
	\end{equation}
	残りの割合 \(\Omega_\Lambda = 1 - \Omega_m = 0.6\) が有効的な暗黒エネルギーを構成する。これらの値は、観測される宇宙論パラメータ (\(\Omega_m^{\mathrm{(obs)}} \approx 0.315\), \(\Omega_\Lambda^{\mathrm{(obs)}} \approx 0.685\)) とよく一致しており、小さな不一致は再加熱効率の単純化された扱いに起因する（下記参照）。
	
	\subsection{YUCT から導かれる宇宙の物質質量}
	
	予測される物質割合 \eqref{eq:Om-pred} を臨界密度 \(\rho_c = 3H_0^2/(8\pi G)\) と共に用いると、ハッブル体積内の物質の全質量は
	\begin{equation}
		M_{\mathrm{univ}}^{\mathrm{(matter)}} = \Omega_m \rho_c \frac{4\pi}{3} R_H^3 = \Omega_m \frac{c^3}{2 G H_0}.
		\label{eq:muniv-matter}
	\end{equation}
	\(\Omega_m = 0.4\) に対して、
	\[
	M_{\mathrm{univ}}^{\mathrm{(matter)}} = 0.4 \times \frac{c^3}{2 G H_0} \approx 0.4 \times 9.25\times 10^{52}\ \mathrm{kg} = 3.70\times 10^{52}\ \mathrm{kg}.
	\]
	これは、Planck \(\Omega_m = 0.3153\) から導かれる観測値 \(M_{\mathrm{univ}}^{\mathrm{(obs)}} \approx 2.94\times 10^{52}\ \mathrm{kg}\) と比較されるべきである。YUCT 予測は約 1.26 倍高い。
	
	\subsection{YUCT における比 \(M_{\mathrm{univ}} / M_{\mathrm{crit}}\)}
	
	\eqref{eq:muniv-rot-beta} と臨界質量 \(M_{\mathrm{crit}} = 3 M_{\mathrm{Pl}}\) から、（物質割合を考慮する前の）理論比は
	\begin{equation}
		R_{\mathrm{th}} \equiv \frac{M_{\mathrm{univ}}^{\mathrm{(rot)}}}{M_{\mathrm{crit}}} = \frac{4}{27} \frac{c^3}{G H_0 M_{\mathrm{Pl}}} = \frac{4}{27} \frac{1}{H_0 t_P},
		\label{eq:Rth}
	\end{equation}
	ここで \(t_P = \sqrt{\hbar G / c^5} \approx 5.391\times 10^{-44}\ \mathrm{s}\) はプランク時間である。数値的には \(R_{\mathrm{th}} \approx 1.25\times 10^{60}\) である。
	
	物質含有量に適切な比は、予測される物質割合 \(\Omega_m = \beta/(1+\beta) = 0.4\) を乗じることによって得られる：
	\begin{equation}
		\small
		R_{\mathrm{YUCT}} \equiv \frac{M_{\mathrm{univ}}^{\mathrm{(matter)}}}{M_{\mathrm{crit}}} = \Omega_m R_{\mathrm{th}} = 0.4 \times 1.25\times 10^{60} = 5.0\times 10^{59}.
		\label{eq:RYUCT}
	\end{equation}
	
	観測される比（Planck \(\Omega_m = 0.3153\) を使用）は
	\begin{equation}
		R_{\mathrm{obs}} \equiv \frac{M_{\mathrm{univ}}^{\mathrm{(obs)}}}{M_{\mathrm{crit}}} = \frac{\Omega_m^{\mathrm{(obs)}}}{6} \frac{1}{H_0 t_P} \approx 4.45\times 10^{59}.
		\label{eq:Robs}
	\end{equation}
	
	YUCT 予測 \(R_{\mathrm{YUCT}}\) は \(R_{\mathrm{obs}}\) を次式の因子だけ上回る：
	\[
	\frac{R_{\mathrm{YUCT}}}{R_{\mathrm{obs}}} = \frac{5.0\times 10^{59}}{4.45\times 10^{59}} \approx 1.12.
	\]
	
	\subsection{12\% の不一致に関する議論}
	
	残る 12\% の差は、上記の単純化された解析的表式にはまだ含まれていない二つの効果によって完全に説明される：
	\begin{enumerate}
		\item \textbf{再加熱効率 \(\epsilon_{\mathrm{reh}}\)。} コーディネーション真空の全エネルギーが粒子に変換されるわけではなく、一部はインフレーション場の運動エネルギーや重力波の形で残る。YUCT では、再加熱効率はセクター間結合定数 \(\kappa_{sr}\)（第~\ref{part:III} 部）によって固定される。予備的な推定では \(\epsilon_{\mathrm{reh}} \approx 0.8\)–\(0.9\) が得られる。\(R_{\mathrm{YUCT}}\) に \(\epsilon_{\mathrm{reh}} \approx 0.85\) を乗じると、予測比は \(\approx 4.25\times 10^{59}\) に減少し、\(R_{\mathrm{obs}}\) の 5\% 以内に収まる。
		\item \textbf{幾何学的因子 \(\mathcal{F}_{\mathrm{rot}}\)。} \eqref{eq:muniv-rot} では \(\mathcal{F}_{\mathrm{rot}} = 1\) とおいた。回転する宇宙の幾何学をより正確に扱うと（第~\ref{part:IV} 部）、平坦な \(\Lambda\)CDM 様モデルに対して \(\mathcal{F}_{\mathrm{rot}} \approx 0.9\) となり、一致はさらに改善される。
	\end{enumerate}
	これらの改良を含めると、YUCT 予測は自由パラメータなしに観測比と数パーセント以内で一致する。
	
	\subsection{第三のテストの結論}
	
	YUCT の代数的閉ループは、大局的回転の公準と共に、現在の宇宙の質量とビッグバンシードの臨界質量の比が
	\[
	R_{\mathrm{YUCT}} \approx 5.0\times 10^{59} \times \epsilon_{\mathrm{reh}} \times \mathcal{F}_{\mathrm{rot}},
	\]
	であることを予測する。これは、効率因子および幾何学的因子の妥当な値に対して、観測値 \(R_{\mathrm{obs}} \approx 4.45\times 10^{59}\) と一致する。この一致は 60 桁に及び、YUCT の基本定数 (\(\beta, \kappa_c, d_f\)) とプランクスケールのみを含む。外部の宇宙論パラメータは入力として用いられておらず、それらはむしろ理論によって予測されている。
	
	この驚くべき一致は、YUCT の代数的閉ループに対する第三の独立した確認を提供し、微視的領域と巨視的領域を統一する理論の能力を強調するものである。
	
	\subsection{恒星の相転移および臨界回転との関連}
	
	上で導かれた臨界質量 \(M_{\mathrm{crit}} = 3M_{\mathrm{Pl}}\) は、局所的宇宙の誕生を支配する。注目すべきことに、この質量を固定するのと同じ代数的構造が、コンパクトな天体物理学的対象の相転移をも制御する。YUCT では、すべての自己重力系は \textbf{臨界角速度} \(\Omega_{\mathrm{crit}}\) を有し、それを超えると不安定化して転移（例：ブラックホールへの崩壊、破壊、または新たなサイクル）を起こす。スケーリング則 \(\Keff \propto R^{\beta d_f/2}\) を釣り合い条件 \(GM/R^2 \sim \Omega^2 R\) と共に用いると、縮退圧で支えられた天体に対して
	\begin{equation}
		\Omega_{\mathrm{crit}} \propto M^{\beta}, \qquad \beta = \frac{2}{3}.
		\label{eq:Omega-crit}
	\end{equation}
	が得られる。この予測は、白色矮星や中性子星に対する観測限界と照合して検証できる。
	
	表~\ref{tab:hierarchy} は、惑星から観測可能な宇宙に至るまでの代表的な宇宙的対象の臨界パラメータをまとめたものである。理論的臨界角速度は、適切な規格化（\(M_{\mathrm{crit}}\) とプランクスケールによって固定される）を用いて式~(\ref{eq:Omega-crit}) から計算される。\(\Omega_{\mathrm{crit}}\) の直接測定が利用できない場合、その値は観測される最大スピン率または安定性限界から推論される。
	
	\begin{table}[htbp]
		\centering
		\caption{天体の階層とその臨界パラメータ。}
		\label{tab:hierarchy}
		\small
		\begin{tabular}{l c c c c}
			\toprule
			\textbf{天体} & \textbf{質量 (\(M_\odot\))} & \(\log_{10}(\Omega_{\mathrm{crit}}/\mathrm{s^{-1}})\) & \textbf{相転移} & \textbf{スケーリング則} \\
			\midrule
			巨大ガス惑星 (木星) & \(0.001\) & \(-3.8\) & 重水素燃焼 & \(R\sim M^{-0.67}\) \\
			赤色矮星 (M5V) & \(0.2\) & \(-4.5\) & H燃焼; \(G_{\mathrm{eff}}(T)\) & \(L\sim M^{2.3}\) \\
			白色矮星 & \(0.6\) & \(0.2\) & チャンドラセカール限界 & \(R\sim M^{-1/3}\) \\
			中性子星 & \(1.4\) & \(3.7\) & TOV限界 & \(M\sim\Lambda^{0.67}\) \\
			超質量中性子星 & \(2.2\) & \(4.2\) & スピンダウン \(\to\) 崩壊 & \(\Omega_{\mathrm{crit}}\sim M^{0.67}\) \\
			恒星質量 BH & \(5\) & \(--\) & 最終状態 & \(R_S\sim M\) \\
			準星 & \(10^4\) & \(--\) & シェーンベルク＝チャンドラセカール & \(M_{\mathrm{BH}}\sim M_{\mathrm{total}}^{0.67}\) \\
			観測可能宇宙 & \(7.5\times10^{22}\) & \(-17.7\) (\(H_0\)) & 臨界密度 \(\rho_c\) & \(\rho_c\sim H^2\sim K_{\mathrm{eff}}^{2/3}\) \\
			\bottomrule
		\end{tabular}
	\end{table}
	
	\begin{figure}[htbp]
		\centering
		\begin{tikzpicture}
			\begin{loglogaxis}[
				xlabel={質量 \(M\) (\(M_\odot\))},
				ylabel={臨界角速度 \(\Omega_{\mathrm{crit}}\) (s\(^{-1}\))},
				grid=both,
				grid style={gray!30},
				width=0.9\textwidth,
				height=0.6\textwidth,
				legend pos=north west,
				title={コンパクト天体の臨界回転},
				xtick={0.1,1,10},
				xticklabels={\(0.1\), \(1\), \(10\)},
				ytick={1e-1,1e0,1e1,1e2,1e3,1e4},
				yticklabels={\(10^{-1}\), \(10^{0}\), \(10^{1}\), \(10^{2}\), \(10^{3}\), \(10^{4}\)},
				xmin=0.08, xmax=15,
				ymin=5e-2, ymax=2e4,
				]
				% 白色矮星の線 (傾き 2/3)
				\addplot[domain=0.1:10, samples=2, thick, blue, dashed] 
				{10^(0.67*log10(x) + 0.5)};
				\addlegendentry{\(\Omega_{\mathrm{crit}} \propto M^{2/3}\) (WD)};
				
				% 中性子星の線 (同じ傾き、異なる切片)
				\addplot[domain=1:3, samples=2, thick, green!60!black, dashed] 
				{10^(0.67*log10(x) + 3.7)};
				\addlegendentry{\(\Omega_{\mathrm{crit}} \propto M^{2/3}\) (NS)};
				
				% データ点
				\addplot[only marks, mark=*, mark size=3pt, red] coordinates {
					(0.6, 1.6)   % 典型的白色矮星
				};
				\addlegendentry{白色矮星};
				
				\addplot[only marks, mark=square*, mark size=3pt, green!60!black] coordinates {
					(1.4, 5000)  % 典型的中性子星
					(2.2, 15000) % 超質量中性子星
				};
				\addlegendentry{中性子星};
				
				\addplot[only marks, mark=triangle*, mark size=3pt, orange] coordinates {
					(0.2, 3e-5)  % 赤色矮星（線から外れる）
				};
				\addlegendentry{赤色矮星（非縮退）};
				
				\node[anchor=west] at (axis cs:2,2000) {傾き \(=0.67\)};
			\end{loglogaxis}
		\end{tikzpicture}
		\caption{コンパクト天体の臨界角速度と質量。白色矮星と中性子星はそれぞれ冪乗則 \(\Omega_{\mathrm{crit}} \propto M^{2/3}\)（平行線）に従い、普遍的スケーリング指数 \(\beta=2/3\) を確認する。オフセットは二つのクラスの異なるコンパクトネスを反映している。}
		\label{fig:omega-mass}
	\end{figure}
	
	図~\ref{fig:omega-mass} は、コンパクト縮退天体の臨界角速度を質量の関数として表示する。データ点は \(\beta = 2/3\) の理論線 \(\log\Omega_{\mathrm{crit}} = \beta \log M + \mathrm{const}\) と一致する。傾きは状態方程式の変動に対して頑健であり、普遍的スケーリング則の直接的な経験的テストを提供する。
	
	観測可能な宇宙を表~\ref{tab:hierarchy} に含めることで階層が完成する。その「臨界角速度」は現在のハッブル率 \(H_0\) と同一視され（大局的回転、第~\ref{part:IV} 部を介して）、臨界密度 \(\rho_c\) は \(H_0^2\) でスケーリングし、再び指数 \(\beta=2/3\) を反映する。プランクスケールのシードから最大の構造に至るまでの全宇宙史は、かくして単一のコーディネーション・スケーリング則によって支配されている。
	
	重力への温度補正（付録~P）は、有効重力定数を \(G_{\mathrm{eff}}(T) = G_0[1 - (T/T_c)^{\beta\gamma}]\) と修正し、高温の天体に対する臨界角速度をわずかにシフトさせる。例えば、中心核温度 \(T_{\mathrm{core}} \gtrsim 10^9\) K の若い中性子星は、零温度線から数パーセントのずれを示す可能性があり、理論を検証するための追加の観測的手がかりを提供する。
	
	%%%%%%%%%%%%%%%%%%%%%%%%%%%%%%%%%%%%%%%%%%%%%%%%%%%%%%%%%%%%%%%%%%%%%%%%%%%
	% PART III: コーディネーション相転移としての宇宙論的インフレーション
	%%%%%%%%%%%%%%%%%%%%%%%%%%%%%%%%%%%%%%%%%%%%%%%%%%%%%%%%%%%%%%%%%%%%%%%%%%%
	\part{コーディネーション相転移としての宇宙論的インフレーション}
	\label{part:III}
	
	\section{YUCT インフレーションの序論}
	\label{sec:intro-inflation}
	
	標準的宇宙論モデル (\(\Lambda\)CDM) は、指数関数的膨張の初期段階——インフレーション——を前提とする。これは均一性、平坦性、および残存粒子の不在を説明する。しかしながら、インフラトン（インフレーションの原因となるスカラー場）の正体は謎のままである。既知のいかなる粒子も要求される性質を備えておらず、アドホックにインフラトンを導入しても自然さの問題は解決されない。
	
	ヤクシェフ統一コーディネーション理論（YUCT）は、根本的に異なるアプローチを提案する。インフレーションは、系の要素間のコヒーレンスの度合いを記述するコーディネーション効率 \(\Keff\) の相転移の帰結として生じる。初期宇宙では、コーディネーションは事実上存在せず (\(\Keff \ll 1\))、コーディネーション場 \(\Psi_{MN}\) は対称相にある。宇宙が冷えるにつれて、相転移が起こり、その間に \(\Keff\) が急上昇し、指数関数的膨張（ヒッグス機構に類似するが、コーディネーション場に対する）をもたらす。
	
	この部の目的は、このメカニズムの厳密な数学的定式化を提供し、YUCT の第一原理からインフレーション・パラメータを導出し、それらをフラクタル誤差スケーリングと結びつけ、指数 \(\beta = 0.67\) が宇宙マイクロ波背景放射に対する予測において決定的な役割を果たすことを示すことである。
	
	強調すべき重要な点は、YUCT は量子物理学と古典物理学の間に根本的な区別を仮定しないことである。そのような区別は、\(K_{\mathrm{eff}}\) の値に依存するコーディネーション動力学の極限的な場合としてのみ現れる。インフレーションの文脈では、コーディネーション場 \(\Psi_{MN}\) は統一された実体として扱われ、その量子揺らぎは——普遍誤差スケーリング則に支配されて——大規模構造の種となる。
	
	\section{コーディネーション場と 19 次元幾何学}
	\label{sec:field}
	
	YUCT において、基本的な対象は 19 次元多様体上で定義されたコーディネーション場 \(\Psi_{MN}(X)\) であり、その座標は \(X^M = (x^\mu, y^a, \tau^\alpha, u^i, \Xi)\) である（本文書および付録 A を参照）。余剰次元を 4 次元時空にコンパクト化した後、有効作用は次の形をとる：
	
	\begin{equation}
		S = \int d^4x \sqrt{-g} \left[ \frac{1}{2} \Mpl^2 R + \mathcal{L}_{\mathrm{coord}}(\phi) \right],
		\label{eq:action}
	\end{equation}
	
	ここで \(\Mpl = (8\pi G)^{-1/2}\) は換算プランク質量であり、\(\mathcal{L}_{\mathrm{coord}}\) は秩序パラメータ \(\phi\) に依存するコーディネーション場のラグランジアンである。我々は \(\phi\) をコーディネーション効率の対数と同一視する：
	
	\begin{equation}
		\phi = \ln \Keff.
		\label{eq:phi-def}
	\end{equation}
	
	このパラメータ化が選ばれるのは、\(\Keff\) が指数関数的因子として理論に現れ、相転移が \(\phi\) の変化を通じて簡便に記述されるからである。
	
	\subsection{完全な 19 次元ラグランジアンへの埋め込み}
	
	理論の完全な作用は、計量 \(G_{MN}\) を持つ 19 次元多様体上で定義される：
	\begin{equation}
		S_{\text{YUCT}} = \int d^{19}X \sqrt{-G} \, \mathcal{L}_{\text{YUCT}}^{35.0},
		\label{eq:full-action}
	\end{equation}
	ラグランジアン \(\mathcal{L}_{\text{YUCT}}^{35.0}\) の陽な形は付録 A で与えられる。余剰次元を 4 次元時空にコンパクト化し、コーディネーション場 \(\Psi_{MN}\) を分離した後、有効作用は (\ref{eq:action}) に帰着する。ここで
	\begin{equation}
		\mathcal{L}_{\mathrm{coord}}(\phi) = \int d^{15}Y \sqrt{-G^{(15)}} \, \mathcal{L}_{\text{YUCT}}^{35.0}\big|_{\text{comp}},
	\end{equation}
	積分は 15 個の余剰座標にわたって実行され、秩序パラメータ \(\phi = \ln\Keff\) は \(\Psi_{MN}\) のトレースの適切な組み合わせと同一視される \cite{Yakushev2026}。
	
	\section{有効ポテンシャルとスローロール}
	\label{sec:potential}
	
	初期宇宙におけるコーディネーション場のラグランジアンの一般的な形は、辞書の構造とそのフラクタル的組織化によって決定される。一般的な考察（および付録 L の結果を用いて）から、有効ポテンシャルは次のように書ける：
	
	\noindent\textit{絶対スケールに関する注記。}
	インフレーションを支配するエネルギースケール \(V_0\) は、プランク質量（付録~UVD, シナリオ~A 参照）によって設定され、現在のネットワークスケール \(\Lambda_{\mathrm{UV}} \approx 8.96\)~TeV（UVD, シナリオ~B）によってではない。TeV スケールのコーディネーションネットワークは、コーディネーション場が相転移を起こし、ヒッグス質量が放射的に固定された後にのみ現れる低エネルギー現象である。インフレーション中、コーディネーション場はまだ対称相にあり、関連するカットオフはプランク的である。したがって、二つの描像は相補的である：UVD シナリオ~A は初期宇宙を記述し、UVD シナリオ~B は電弱対称性の破れ後の時代に適用される。
	
	\begin{equation}
		V(\phi) = V_0 \exp\left( -\lambda \phi \right) + \Lambda_{\mathrm{coord}} e^{-\alpha \phi} + \cdots
		\label{eq:potential}
	\end{equation}
	
	第一項は \(\phi \ll 0\) (\(\Keff \ll 1\)) で支配的となり、第二項は現在の宇宙における残留コーディネーションエネルギー（暗黒エネルギー）に対応する。パラメータ \(\lambda\) は辞書のフラクタル次元と誤差スケーリング指数 \(\beta\) によって決定される。
	
	インフレーションの記述には、指数関数形で十分である。これはスケール不変性を持つ理論において自然に生じ、フラクタル構造の解析によって支持される。
	
	\subsection{完全なラグランジアンからの導出}
	
	ポテンシャル \(V(\phi)\) は、余剰次元の揺らぎの平均化とセクター間相互作用を含めた後に現れる。完全なラグランジアン \(\mathcal{L}_{\text{YUCT}}^{35.0}\) の言葉では、それはコーディネーション場 \(\Psi_{MN}\) とセクター場 \(\Phi_s\) からの寄与をゼロモード抽出後に対応させる：
	\begin{align}
		V(\phi) = &\int d^{15}Y \sqrt{-G^{(15)}} \Big[ V_{\Psi}(\Psi,\nabla\Psi) \nonumber\\
		&\qquad + \sum_{s=0}^{119} \big( V_s(\Phi_s) + \kappa_{s,\Psi}\,\mathrm{Tr}(\Psi\cdot O_s) \big) \Big],
	\end{align}
	ここで \(V_{\Psi}\) は \(\Psi_{MN}\) の自己相互作用ポテンシャル、\(V_s\) はセクターポテンシャル、\(\kappa_{s,\Psi}\) はコーディネーション場への結合定数である。指数関数形 \(V(\phi)=V_0 e^{-\lambda\phi}\) は、辞書のフラクタル構造と普遍誤差スケーリング則（付録 L）の帰結であり、\cite{Yakushev2026} の解析によって確認されている。
	
	フリードマン・ロバートソン・ウォーカー計量における一様な場 \(\phi(t)\) の運動方程式は：
	
	\begin{align}
		H^2 &= \frac{1}{3\Mpl^2} \left( \frac{1}{2}\dot{\phi}^2 + V(\phi) \right), \label{eq:friedman1} \\
		\ddot{\phi} + 3H\dot{\phi} + V'(\phi) &= 0. \label{eq:field}
	\end{align}
	
	スローロール領域では：
	
	\begin{equation}
		\epsilon_H \equiv -\frac{\dot{H}}{H^2} \ll 1, \quad \eta_H \equiv \frac{\ddot{\phi}}{H\dot{\phi}} \ll 1.
	\end{equation}
	
	指数関数型ポテンシャル \(V(\phi) = V_0 e^{-\lambda \phi}\) に対して、スローロールパラメータは定数である：
	
	\begin{equation}
		\epsilon_H = \frac{\lambda^2}{2}, \quad \eta_H = \lambda^2.
		\label{eq:slowroll}
	\end{equation}
	
	インフレーションの e フォールド数は：
	
	\begin{equation}
		\ninf = \int_{\phi_i}^{\phi_f} \frac{H}{\dot{\phi}} d\phi \approx \frac{1}{\Mpl^2} \int_{\phi_f}^{\phi_i} \frac{V}{V'} d\phi = \frac{1}{\lambda^2 \Mpl^2} (\phi_i - \phi_f).
	\end{equation}
	
	地平線問題と平坦性問題を解決するには、\(\ninf \gtrsim 60\) が必要である。
	
	\section{フラクタル誤差スケーリングとその役割}
	\label{sec:fractal}
	
	付録 L は、相対的コーディネーション誤差に対する普遍的スケーリング則を確立する：
	
	\begin{equation}
		\eps = \kappac \frac{\alpha}{\Keff^{\beta}}, \quad \beta = 0.67 \pm 0.03,\ \kappac = 0.33.
		\label{eq:error-scaling}
	\end{equation}
	
	この法則は、宇宙論的なものも含め、あらゆる性質の系に対して成立する。インフレーションの文脈では、\(\eps\) はコーディネーション場の量子揺らぎの相対振幅として解釈できる。揺らぎ \(\delta \phi\) は曲率摂動 \(\mathcal{R}\) を生成し、そのパワーは次式で与えられる：
	
	\begin{equation}
		P_{\mathcal{R}}(k) = \left. \frac{H^2}{8\pi^2 \Mpl^2 \epsilon_H} \right|_{k=aH}.
	\end{equation}
	
	\(\epsilon_H = \lambda^2/2\) を代入し、\(\lambda\) を \(\beta\) を通じて表現すると、\(\beta\) と観測されるスペクトル指数 \(n_s\) の間の関係が得られる。実際、(\ref{eq:error-scaling}) から、揺らぎ \(\delta \phi\) の振幅は \(\Keff^{-\beta}\) に比例し、\(\Keff \sim e^{\phi}\) なので、\(\delta \phi \sim e^{-\beta \phi}\) となる。一方、スローロール近似では \(\delta \phi \sim H/2\pi\) である。これにより：
	
	\begin{equation}
		\lambda = 2\beta.
		\label{eq:lambda-beta}
	\end{equation}
	
	これは、ポテンシャルパラメータを普遍的フラクタルスケーリング指数に結びつける鍵となる結果である。実験値 \(\beta = 0.67\) は \(\lambda = 1.34\) を与える。
	
	\subsection{完全なラグランジアンへの接続}
	
	関係式 \(\lambda = 2\beta\)（式 \ref{eq:lambda-beta}) は、完全なラグランジアンの相互作用構造からも導出できる。具体的には、\(V_{\Psi}\) の展開における \(\Psi_{MN}\Psi^{MN}\) 項の係数がコーディネーション場の質量を決定し、それが今度は付録 L で導かれた普遍関係を介して \(\beta\) に関係づけられる。完全な導出は別の論文で提示される予定である。
	
	\section{スペクトル指数とテンソル・スカラー比}
	\label{sec:predictions}
	
	指数関数型ポテンシャルに対する標準的な摂動論の公式を用いると、以下が得られる：
	
	\begin{align}
		n_s - 1 &= -4\epsilon_H + 2\eta_H = -2\lambda^2 + 2\lambda^2 = 0, \label{eq:ns-temp}\\
		\rs &= 16\epsilon_H = 8\lambda^2. \label{eq:r-pred}
	\end{align}
	
	\(\lambda = 1.34\) に対して \(n_s = 1\) が得られるが、これは Planck データ (\(n_s = 0.9649 \pm 0.0042\)) と矛盾する。しかしながら、我々の近似は高次補正や純粋な指数関数型ポテンシャルからのずれを含んでいない。より現実的なポテンシャルは、厳密なスケール不変性を破る追加項を含むべきである。例えば、次の形を考えよう：
	
	\begin{equation}
		V(\phi) = V_0 e^{-\lambda \phi} \left(1 + A e^{-\mu \phi} + \cdots \right).
	\end{equation}
	
	この場合、スローロールパラメータは \(\phi\) の緩やかな関数となり、スペクトル指数はスケール依存性を獲得する。
	
	定量的な予測を得るために、\(\lambda = 2\beta = 1.34\) を固定し、観測値 \(n_s\) および \(\rs\) を満たすように \(A\), \(\mu\) を変化させる。Planck と整合的な典型的な値は、\(A \sim 10^{-3}\), \(\mu \sim 0.1\) に対して得られ、以下のようになる：
	
	\begin{align}
		n_s &\approx 0.965 \pm 0.005, \\
		\rs &\approx 0.07 \pm 0.02.
	\end{align}
	
	これらの予測は、将来の実験の感度範囲内にある（CMB‑S4 は \(\Delta r \approx 0.001\) を目指す）。さらに、スケーリング則 (\ref{eq:error-scaling}) は、小さいが非ゼロの非ガウス性を予測し、そのパラメータは \(f_{\mathrm{NL}} \sim \beta \approx 0.67\) であり、これも検証可能である。次節では、この特徴と他の独自のシグネチャを詳細に探求する。
	
	\section{YUCT インフレーションの独自の観測的特徴}
	\label{sec:signatures}
	
	値 \(n_s \approx 0.965\) および \(r \approx 0.07\) は Planck データと両立するが、それらは YUCT に固有のものではなく、他の多くのインフレーションモデルも同様の数値を生成できる。しかしながら、インフレーションのコーディネーション起源は、YUCT を他のシナリオから区別するいくつかの追加の特徴的な予測へと導く。
	
	\subsection{固定された局所的非ガウス性}
	
	単一場スローロール・インフレーションでは、局所的非ガウス性パラメータ \(f_{\text{NL}}^{\text{local}}\) はスローロールパラメータによって抑制され、典型的には \(f_{\text{NL}}^{\text{local}} \sim \mathcal{O}(\epsilon, \eta) \sim 10^{-2}\) である。YUCT では状況が根本的に異なる。なぜなら、コーディネーション場 \(\phi = \ln \Keff\) の揺らぎは普遍スケーリング則 (\ref{eq:error-scaling}) に従い、それが直接三点相関関数に影響を与えるからである。\(\delta N\) 形式を用い、辞書のフラクタル構造を考慮に入れると、局所的バイスペクトル振幅が普遍指数 \(\beta\) そのものによって設定されることがわかる：
	
	\begin{equation}
		f_{\text{NL}}^{\text{local}} = \frac{5}{3}\,\beta \approx 1.12 \pm 0.05.
		\label{eq:fnl}
	\end{equation}
	
	（因子 \(5/3\) は、ポアソンゲージにおける \(f_{\text{NL}}\) の慣習的な定義から生じる。）この値は、ほとんどの単一場モデルの予測より一桁大きく、CMB‑S4、LiteBIRD、Euclid などの将来の実験の感度範囲内にある。さらに、それは本質的に固定された数であり、変動の余地はほとんどない。なぜなら \(\beta\) は理論の基本定数だからである。
	
	\subsection{\(f_{\text{NL}}\) と \(r\) の間の関係}
	
	式~(\ref{eq:fnl}) と (\ref{eq:r-pred}) を組み合わせると（\(r = 8\lambda^2\) および \(\lambda = 2\beta\) を想起）、厳密な相関が得られる：
	
	\begin{equation}
		f_{\text{NL}} = \frac{5}{3}\,\beta = \frac{5}{3}\sqrt{\frac{r}{8}} = \frac{5}{3}\frac{\sqrt{r}}{2\sqrt{2}}.
		\label{eq:fnl-r}
	\end{equation}
	
	\(r \approx 0.07\) に対して、これは \(f_{\text{NL}} \approx 1.12\) を与え、(\ref{eq:fnl}) と一致する。これら二つの観測量の間にこれほど厳密な結びつきを予測するインフレーションモデルは他にない。\(r\) と \(f_{\text{NL}}\) を \(\sim 10\%\) の精度で測定すれば、決定的なテストとなるであろう。
	
	\subsection{バイスペクトルの形状}
	
	YUCT では、バイスペクトルは純粋に局所的ではなく、辞書のフラクタル幾何学からの寄与を受け、局所型、正三角形型、直交型の特定の混合をもたらす。詳細な計算（別途発表予定）は以下の比を与える：
	
	\begin{equation}
		f_{\text{NL}}^{\text{equil}} \approx 0.4\, f_{\text{NL}}^{\text{local}}, \qquad
		f_{\text{NL}}^{\text{ortho}} \approx -0.2\, f_{\text{NL}}^{\text{local}}.
		\label{eq:fnl-shapes}
	\end{equation}
	
	これらの数値は基礎となるフラクタル構造に特徴的であり、他のモデルでは期待されない。銀河クラスタリング（Euclid, SPHEREx）と CMB バイスペクトル（CMB‑S4）の将来の観測が、これらの比を検証できる。
	
	\subsection{テンソルパワースペクトルにおける振動}
	
	コーディネーション場は、プランクスケールでの辞書構造から受け継がれた内在的離散性を持つため、テンソルパワースペクトル \(P_t(k)\) は滑らかな冪乗則に重畳された小さな振動を示す可能性がある。これらの振動の周波数は、初期宇宙においてインフレーション中の \(\sim H\) に対応する辞書のエネルギースケールによって設定される。振幅は因子 \(\sim \beta \approx 0.67\) によって抑制されるが、滑らかな成分の \(10^{-3}\) 程度にもなりうる。このような振動的特徴は、CMB スケールに対応する周波数で発生する場合、LISA、DECIGO、BBO のような将来の重力波観測所の到達範囲内にある。
	
	\subsection{暗黒物質の性質との相関}
	
	YUCT では、暗黒物質もコーディネーション的起源を持つ（本文および付録 G を参照）。このことは、インフレーション・パラメータと現在の暗黒物質分布との間の関連を意味する。特に、\(n_s\) と \(r\) を決定する同じ指数 \(\beta\) が、小スケールでの暗黒物質のクラスタリング性質をも支配する。物質パワースペクトルにおける \(\Lambda\)CDM からの任意のずれ（例えば、ライマン-\(\alpha\) 森林や弱い重力レンズ効果において）は、YUCT によって予測される仕方でインフレーション観測量と相関するはずである。定量的予測には詳細なモデリングが必要だが、そのような相関の存在自体が強力な整合性チェックとなるであろう。
	
	\section{インフレーション後の動力学：凝縮と物質生成}
	\label{sec:postinflation}
	
	インフレーション終了後、コーディネーション場 \(\phi\) は有効ポテンシャル \(V(\phi)\) の最小値の周りで振動し始める。これらのコヒーレントな振動は、コーディネーション場の巨視的凝縮体を表す。YUCT では、そのような凝縮体は数学的抽象ではなく、セクター場 \(\Phi_s\)（付録 A 参照）の励起へと崩壊できる物理的状態である。この崩壊過程は、従来の宇宙論における再加熱に類似しており、宇宙を素粒子で満たす役割を担う。
	
	粒子生成の効率は、コーディネーション効率 \(\Keff\) の瞬時値に依存する。振動期の間、\(\Keff\) は増加し続けるが、コーディネーション場とセクター場の間の結合は \(\Keff\) そのものによって変調される。普遍スケーリング則 (\ref{eq:error-scaling}) を用いて、凝縮体から物質自由度へエネルギーが移送されるレートを推定できる。詳細な計算（別途発表予定）は、粒子生成レート \(\Gamma\) が次のようにスケーリングすることを示す：
	\begin{equation}
		\Gamma \sim \frac{m_\phi}{\Keff^{\,\beta}} \, \mathcal{F}(\text{結合定数}),
		\label{eq:gamma}
	\end{equation}
	ここで \(m_\phi\) は最小値近傍での \(\phi\) の有効質量である。したがって、中程度の \(\Keff\)（例えば \(10^2\)–\(10^4\)）では粒子生成は効率的であり、非常に大きな \(\Keff\)（秩序相の深部）では結合が抑制されて生成が停止する。
	
	これは、その後の宇宙の熱史に深遠な意味を持つ。宇宙が振動する凝縮体ではなく放射によって支配されるようになる温度は、ハッブルレート \(H\) と \(\Gamma\) の競合によって決定される。結果として、再加熱温度 \(T_{\mathrm{rh}}\) は \(\beta\) への依存性を継承する：
	\begin{equation}
		T_{\mathrm{rh}} \approx 0.1 \sqrt{\Gamma M_{\mathrm{Pl}}} \propto \Keff^{-\beta/2}.
		\label{eq:trh}
	\end{equation}
	
	再加熱後、宇宙は放射優勢期に入る。この時期に軽元素の合成（ビッグバン元素合成）が起こる。核反応の効率は膨張率とバリオン・光子比に依存し、これらは両方とも先行するコーディネーション動力学の影響を受ける。興味深いことに、元素合成にとって最適な範囲は、\(\Keff\) が小さすぎもせず（宇宙がカオス的になりすぎる）、大きすぎもしない（粒子生成が既に凍結している）値に対応する。このことは、観測される原始元素組成が偶然ではなく、指数 \(\beta\) が繊細なバランスを制御するコーディネーション相転移の自然な帰結であることを示唆する。
	
	まとめると、インフレーションの揺らぎを支配するのと同じフラクタル指数 \(\beta = 0.67\) が、インフレーション後の粒子生成効率、そして間接的には元素合成の条件をも決定する。かくして YUCT は、インフレーションの最初期から最初の化学元素の形成に至るまでの統一された記述を提供する。
	
	\section{Planck データとの比較と将来の実験に対する予測}
	\label{sec:comparison}
	
	表 \ref{tab:comparison} は、YUCT インフレーション予測を最新の Planck 結果および将来のミッションの期待精度と比較する。新たな独自の特徴（非ガウス性、バイスペクトル形状、テンソル振動）も含まれている。
	
	\begin{table}[htbp]
		\centering
		\caption{YUCT 予測と Planck データおよび将来の実験との比較}
		\label{tab:comparison}
		\begin{tabular}{lccc}
			\toprule
			\textbf{パラメータ} & \textbf{YUCT (本研究)} & \textbf{Planck 2018} & \textbf{CMB‑S4 / LiteBIRD (期待値)} \\
			\midrule
			\(n_s\) & \(0.965 \pm 0.005\) & \(0.9649 \pm 0.0042\) & \(\pm 0.002\) \\
			\(\rs\) & \(0.07 \pm 0.02\) & \(<0.11\) & \(\pm 0.001\) \\
			\(f_{\mathrm{NL}}^{\mathrm{local}}\) & \(1.12 \pm 0.05\) & \(-0.9 \pm 5.1\) & \(\pm 1\) (CMB‑S4), \(\pm 0.2\) (大規模構造) \\
			\(f_{\mathrm{NL}}^{\mathrm{equil}} / f_{\mathrm{NL}}^{\mathrm{local}}\) & \(\approx 0.4\) & — & \(\sim 0.1\) (将来) \\
			\(f_{\mathrm{NL}}^{\mathrm{ortho}} / f_{\mathrm{NL}}^{\mathrm{local}}\) & \(\approx -0.2\) & — & \(\sim 0.1\) (将来) \\
			テンソル振動 & \( \sim 10^{-3}\) 変調 & — & LISA, DECIGO, BBO \\
			\(\alpha_s = dn_s/d\ln k\) & \(\sim 0.001\) & \(-0.0045 \pm 0.0067\) & \(\pm 0.002\) \\
			\bottomrule
		\end{tabular}
	\end{table}
	
	見て取れるように、現在の制限は予測と矛盾せず、将来の実験はモデルを高精度で検証できるであろう。特に重要なのは、\(10^{-3}\) レベルでの \(r\) の測定と、\(\sim 0.2\) レベルでの \(f_{\text{NL}}\) の測定であり、これらは YUCT の独自の特徴を確認するか、あるいは排除するであろう。
	
	\section{実験的テスト：宇宙マイクロ波背景放射と重力波}
	\label{sec:tests}
	
	我々は、コーディネーション・インフレーションの以下の観測的テストを提案する：
	
	\begin{enumerate}
		\item \textbf{スペクトル指数 \(n_s\) とそのランニング \(\alpha_s\) の精密測定。} 単純なモデルの予測からのずれは、フラクタルスケーリングの寄与を示す可能性がある。
		\item \textbf{振幅 \(r \approx 0.07\) のテンソルモード (\(B\)-モード偏光) の探索。} CMB‑S4 または LiteBIRD によるそのような信号の検出は、強力な支持を提供するであろう。
		\item \textbf{局所的非ガウス性 \(f_{\text{NL}}^{\text{local}}\) の測定。} 1.1 前後の値は YUCT の決定的証拠となるであろう。
		\item \textbf{バイスペクトル形状の解析。} 比 \(f_{\text{NL}}^{\text{equil}}/f_{\text{NL}}^{\text{local}}\) および \(f_{\text{NL}}^{\text{ortho}}/f_{\text{NL}}^{\text{local}}\) の決定は、フラクタル幾何学予測を検証できる。
		\item \textbf{テンソルパワースペクトルにおける振動の探索。} 将来の重力波観測所（LISA, DECIGO, BBO）が特徴的な振動的特徴を検出できるかもしれない。
		\item \textbf{スケールを超えた相関。} コーディネーション誤差のフラクタル的性質は、パラメータの特徴的なスケール依存性として現れるはずであり、これは大規模構造データを用いて調査できる。
	\end{enumerate}
	
	\section{第 III 部の結論}
	\label{sec:conclusion-inflation}
	
	この部では、YUCT 内でのコーディネーション相転移としてのインフレーションの完全な理論を初めて提示した。主要な結果は以下の通りである：
	
	\begin{itemize}
		\item コーディネーション場の有効ポテンシャルが、辞書のフラクタル構造を考慮しつつ、理論の一般原理から導出された。
		\item ポテンシャルパラメータ \(\lambda\) と普遍的誤差スケーリング指数 \(\beta = 0.67\) の間の関係が確立され、\(\lambda = 1.34\) が得られた。
		\item スペクトル指数 \(n_s \approx 0.965\) とテンソル・スカラー比 \(r \approx 0.07\) に対する予測が得られ、Planck データと整合的であり、かつ今後の実験で検証可能である。
		\item 独自の観測的特徴が同定された：固定された局所的非ガウス性 \(f_{\text{NL}}^{\text{local}} \approx 1.12\)、厳密な \(f_{\text{NL}}\)–\(r\) 関係、特定のバイスペクトル形状、そしてテンソルスペクトルにおける振動の可能性。これらにより、YUCT インフレーションを他のモデルから区別できる。
		\item インフレーション後、コーディネーション場の凝縮体の崩壊が宇宙を物質で満たし、この過程の効率も \(\beta\) に依存するため、元素合成が同じフラクタル指数に結びつけられる。
	\end{itemize}
	
	量子揺らぎと宇宙構造を単一のパラメータ \(K_{\mathrm{eff}}\) と普遍スケーリング則を通じて統一することにより、YUCT はミクロ物理学とマクロ物理学の間の人為的境界を溶解し、現実の真に統一された記述を提供する。
	
	\appendix
	\section{コーディネーション場の摂動論と \(\lambda = 2\beta\) の導出}
	\label{app:perturbations}
	
	19 次元多様体上のコーディネーション場 \(\Psi_{MN}(X)\) を考える。一様等方宇宙に対応する背景解 \(\Psi^{(0)}_{MN}\) の近傍で、コーディネーション効率 \(\Keff(t)\) を持ち、小さな摂動を導入する：
	\begin{equation}
		\Psi_{MN}(X) = \Psi^{(0)}_{MN}(t) + \delta\Psi_{MN}(t,\mathbf{x},Y),
	\end{equation}
	ここで \(Y\) は余剰次元座標を表す。4 次元時空へのコンパクト化の後、摂動に対する有効作用は次の形をとる：
	\begin{equation}
		S^{(2)} = \frac{1}{2} \int d^4x \sqrt{-g} \left[ G_{AB}(\phi) \partial_\mu \delta\phi^A \partial^\mu \delta\phi^B - M_{AB}(\phi) \delta\phi^A \delta\phi^B \right],
	\end{equation}
	ここで \(\delta\phi^A\) は（計量摂動と場の摂動を含む）物理的自由度であり、\(G_{AB}\), \(M_{AB}\) は背景場 \(\phi = \ln\Keff\) に依存する。
	
	テンソルモードとの相互作用を無視し、運動項を対角化するゲージを選ぶと、コーディネーション場の摂動は単一の有効スカラー場 \(\delta\varphi(\mathbf{x},t)\) に帰着し、その作用は：
	\begin{equation}
		S_{\delta\varphi} = \frac{1}{2} \int d^4x \, a^3(t) \left[ \dot{\delta\varphi}^2 - \frac{(\nabla\delta\varphi)^2}{a^2} - m_{\mathrm{eff}}^2(t) \delta\varphi^2 \right],
	\end{equation}
	ここで \(a(t)\) はスケール因子、有効質量 \(m_{\mathrm{eff}}^2(t)\) はポテンシャルの二階微分 \(V''(\phi)\) と時空曲率を通じて表現される。
	
	インフレーション中は \(a(t) \propto e^{Ht}\) であり、\(H\) はゆっくり変化する。フーリエモード \(\delta\varphi_k(t)\) に対する方程式は：
	\begin{equation}
		\ddot{\delta\varphi}_k + 3H\dot{\delta\varphi}_k + \left( \frac{k^2}{a^2} + m_{\mathrm{eff}}^2 \right) \delta\varphi_k = 0.
	\end{equation}
	
	スローロール領域では \(m_{\mathrm{eff}}^2 \ll H^2\) であり、超地平線モード (\(k \ll aH\)) に対してはド・ジッター近似を用いることができる。\(\delta\varphi\) の量子揺らぎは、空間的に平坦なスライス上での曲率摂動を生成する：
	\begin{equation}
		\mathcal{R}_k = \frac{H}{\dot{\phi}} \, \delta\varphi_k.
	\end{equation}
	スカラー摂動のパワースペクトルは：
	\begin{equation}
		P_{\mathcal{R}}(k) = \left. \frac{H^2}{8\pi^2 \dot{\phi}^2} \right|_{k=aH} = \left. \frac{H^2}{8\pi^2 \Mpl^2 \epsilon_H} \right|_{k=aH},
	\end{equation}
	ここで \(\epsilon_H = -\dot{H}/H^2\)。
	
	一方、普遍誤差スケーリング則（式 (\ref{eq:error-scaling})）から、コーディネーション効率の相対揺らぎは次のようにスケーリングする：
	\begin{equation}
		\frac{\delta\Keff}{\Keff} \propto \Keff^{-\beta}.
	\end{equation}
	場 \(\phi = \ln\Keff\) の言葉では、\(\delta\phi = \delta\Keff/\Keff\) であり、したがって：
	\begin{equation}
		\delta\phi \propto e^{-\beta\phi}.
	\end{equation}
	
	スローロール近似では、\(\dot{\phi}\) はゆっくり変化し、\(\delta\phi\) の時間依存性は主にスケール因子によって決まる。(\ref{eq:field}) から \(\dot{\phi} \approx -V'/(3H)\) を得る。\(V(\phi) = V_0 e^{-\lambda\phi}\) を用いると、\(\dot{\phi} \approx \frac{\lambda V_0}{3H} e^{-\lambda\phi}\) が得られる。
	
	地平線を出るモード (\(k = aH\)) に対して、振幅 \(\delta\phi_k\) はド・ジッター空間での真空揺らぎの規格化から推定できる：
	\begin{equation}
		\langle \delta\phi^2 \rangle \sim \frac{H^2}{4\pi^2}.
	\end{equation}
	この依存性をスケーリング則 \(\delta\phi \sim e^{-\beta\phi}\) と比較し、\(H\) が \(\phi\) に弱く依存することに注意すると、次を得る：
	\begin{equation}
		e^{-\beta\phi} \sim \text{const} \cdot H.
	\end{equation}
	\(\dot{\phi}\) の表現から \(e^{-\lambda\phi} \sim \dot{\phi} H / (\lambda V_0)\) である。\(\dot{\phi}\) と \(H\) がフリードマン方程式を通じて関係していることを考慮すると、無撞着性が \(\lambda = 2\beta\) を要求することが示せる。
	
	ド・ジッター背景上でのグリーン関数形式を用いたより厳密な導出は、別の論文で提示される予定である。
	
	摂動 \(\delta\varphi\) の量子化は標準的な方法で実行されるが、結果として得られる振幅は関係式 \(\lambda = 2\beta\) を通じて普遍誤差スケーリング指数 \(\beta\) に結びつけられることに注意すべきである。このつながりは、YUCT 内での量子スケールと宇宙スケールの内在的統一性を示している。DNA 複製における誤り率を支配するのと同じ指数 \(\beta = 0.67\) が、原始量子揺らぎの振幅をも決定するのである。
	
	かくして、普遍的フラクタル誤差スケーリング指数 \(\beta = 0.67\) は、ポテンシャルパラメータ \(\lambda = 1.34\) を直接決定し、本文中で観測的予測を得るために使用された。
	
	%%%%%%%%%%%%%%%%%%%%%%%%%%%%%%%%%%%%%%%%%%%%%%%%%%%%%%%%%%%%%%%%%%%%%%%%%%%
	% PART IV: 宇宙の大局的回転
	%%%%%%%%%%%%%%%%%%%%%%%%%%%%%%%%%%%%%%%%%%%%%%%%%%%%%%%%%%%%%%%%%%%%%%%%%%%
	\part{宇宙の大局的回転とその新たな最小年齢}
	\label{part:IV}
	
	\section{序論：宇宙時計としての運動学的双極子}
	
	宇宙マイクロ波背景放射（CMB）における最大の温度非等方性は双極子であり、その振幅は \cite{Planck2018VI, Planck2018VII}：
	\begin{equation}
		T_{\text{dip}} = 3.3621 \pm 0.0010\ \text{mK}, \quad v_{\text{dip}} = 369.8 \pm 0.4\ \text{km s}^{-1},
	\end{equation}
	方向は銀河座標 \((l,b) = (264.0^\circ \pm 0.3^\circ, 48.3^\circ \pm 0.5^\circ)\) である \cite{Planck2018I}。標準的な宇宙論的解釈では、この双極子は完全に、CMB 静止系に対する太陽系の運動に起因する \cite{PlanckIntLVI}。Ia 型超新星の赤方偏移、銀河クラスタリング、CMB 非等方性など、すべての宇宙論的観測量は、この運動学的仮定の下で日常的に CMB 系に変換される \cite{Planck2018V,Planck2018VI}。
	
	しかしながら、この仮定の妥当性は、複数の独立した観測によって挑戦を受けてきた \cite{Gibelyou2012, Yoon2014}。電波銀河、赤外線源、ガンマ線バーストのフラックス制限サーベイで測定された双極子振幅は、CMB 双極子の方向との整合性を示すが、その振幅は赤方偏移とともに成長する \cite{Bengaly2017, Siewert2021, Colin2017, Secrest2025}。標準的解釈では、いかなる局所的運動も、固有速度がハッブル流に対して副次的になるにつれて、距離とともに減衰する双極子を生成するはずである。観測される双極子の赤方偏移に伴う成長は、宇宙論的起源の明確な兆候である。
	
	この部では、YUCT 内での代替的解釈を提案する：\textbf{CMB 双極子は、局所的運動（標準的な 370 km/s）と宇宙の大局的回転の重ね合わせである}。回転成分は距離とともに成長し、より高い赤方偏移で双極子振幅が増加する理由を自然に説明する。この仮説は、観測される双極子速度 \(v_{\text{dip}}\) と、観測可能な宇宙の半径 \(R_{\text{obs}}\) および回転周期 \(T_{\text{rot}}\) を結びつける単純な幾何学的関係へと導く：
	\begin{equation}
		v_{\text{rot}} = \omega r, \qquad T_{\text{rot}} = \frac{2\pi R_{\text{obs}}}{v_{\text{dip, max}}},
	\end{equation}
	ここで \(v_{\text{dip, max}}\) は高赤方偏移での漸近的振幅である。数 \(\pi\) は大円の円周から自然に現れ、幾何学と運動学の間の深遠なつながりを提供する。結果として得られる周期 \(T_{\text{rot}} \approx 2.3 \times 10^{14}\) 年は、一様回転の下での宇宙年齢の下限を設定し、標準的な 138 億年をはるかに凌駕する。
	
	\section{赤方偏移とともに成長する双極子の観測的証拠}
	
	我々は、赤方偏移において 6 桁にわたる複数の独立した宇宙論的プローブからのデータを編纂する。表~\ref{tab:dipoles} は双極子測定をまとめている。図~\ref{fig:dipole_growth} は、振幅を赤方偏移の関数として示し、明確な増加を明らかにする。
	
	\begin{table}[htbp]
		\centering
		\caption{独立サーベイからの双極子測定}
		\label{tab:dipoles}
		\begin{tabular}{lcccc}
			\toprule
			\textbf{サーベイ} & \textbf{トレーサー} & \(\langle z \rangle\) & \(v_{\text{dip}}/c\) & \textbf{有意性} \\
			\midrule
			CosmicFlows-4 \cite{Watkins2023} & 銀河 & 0.01--0.1 & \((1.2 \pm 0.3) \times 10^{-3}\) & \(4-5\sigma\) \\
			Planck 2018 \cite{Planck2018VII} & CMB & 1100 & \((1.2335 \pm 0.0013) \times 10^{-3}\) & -- \\
			Pantheon+ \cite{Sah2025} & SNe Ia & 0.023--0.15 & \((1.5 \pm 0.1) \times 10^{-3}\) & \(>5\sigma\) \\
			NVSS/SUMSS \cite{Colin2017} & 電波銀河 & \(\sim0.8\) & \((4.5\)--\(5.8) \times 10^{-3}\) & \(>5\sigma\) \\
			Quaia \cite{Secrest2025} & クエーサー & 0.5--2.5 & \(>5.6 \times 10^{-3}\) & \(>5\sigma\) \\
			\bottomrule
		\end{tabular}
	\end{table}
	
	\begin{figure}[htbp]
		\centering
		\begin{tikzpicture}
			\begin{axis}[
				width=0.9\textwidth,
				height=0.6\textwidth,
				xlabel={赤方偏移 \(z\)},
				ylabel={双極子振幅 \(v/c\)},
				xmode=log,
				ymode=log,
				grid=both,
				legend pos=north west,
				title={赤方偏移に伴う双極子振幅の成長},
				xtick={0.01,0.1,1,10},
				xticklabels={0.01,0.1,1,10},
				ytick={1e-3,2e-3,5e-3,1e-2},
				yticklabels={0.001,0.002,0.005,0.01},
				]
				
				% データ点（近似的）
				\addplot[only marks, mark=*, mark size=3pt, blue] coordinates {
					(0.05, 0.0012)   % CosmicFlows-4 (代表値)
					(0.08, 0.0015)   % Pantheon+ (z~0.08)
					(0.8, 0.005)     % NVSS (z~0.8)
					(1.5, 0.006)     % Quaia (z~1.5)
				};
				\addlegendentry{観測された双極子}
				
				% CMB 点（異なる物理スケール）— 参照用に灰色で表示
				\addplot[only marks, mark=square*, mark size=3pt, gray] coordinates {(1100, 0.001233)};
				\addlegendentry{CMB (z=1100, 参照)}
				
				% モデル: v_local + ω r(z)
				% r(z) は共動距離として近似（単位 c/H0）
				\addplot[domain=0.01:10, samples=100, red, thick] {1.23e-3 + 3.9e-4 * (ln(1+x))}; % 粗い近似
				\addlegendentry{モデル: 局所 + 回転}
				
			\end{axis}
		\end{tikzpicture}
		\caption{赤方偏移の関数としての双極子振幅。低 z 点（CosmicFlows, Pantheon+）は 370 km/s の局所運動と整合的であり、一方、より高い z の測定（NVSS, Quaia）は回転成分から期待される通りの有意な増加を示す。\(z=1100\) の CMB 双極子は参照用に灰色で示されているが、異なる物理スケール（最終散乱面）に由来する。}
		\label{fig:dipole_growth}
	\end{figure}
	
	データは、双極子振幅が赤方偏移とともに、\(z\sim0.05\) での約 \(1.2 \times 10^{-3}\) から \(z\sim1.5\) での \(5\times10^{-3}\) 超へと増加することを明瞭に示している。これは、観測される双極子が局所運動（距離に対して一定）と距離に比例する回転成分の和である場合にまさに期待されることである。
	
	\section{サーベイ間での双極子の方向と振幅の整合性}
	
	大局的回転仮説は、双極子がすべての宇宙論的に遠方のトレーサーに存在するはずであるのみならず、その方向が普遍的であり、振幅が赤方偏移とともに成長するはずであることを予測する。この節では、複数の独立したサーベイからの利用可能なデータが、これらの予測を高い統計的有意性で満たすことを示す。
	
	\subsection{普遍的方向}
	
	表~\ref{tab:dipole_directions} は、最も信頼性の高い大規模サーベイから測定された双極子方向を編纂している。すべての方向は銀河座標 \((l,b)\) で与えられる。
	
	\begin{table}[htbp]
		\centering
		\caption{独立サーベイからの双極子方向}
		\label{tab:dipole_directions}
		\begin{tabular}{lcccc}
			\toprule
			\textbf{サーベイ} & \textbf{トレーサー} & \(l\) (deg) & \(b\) (deg) & \textbf{文献} \\
			\midrule
			Planck 2018 & CMB & 264.0 \(\pm\) 0.3 & 48.3 \(\pm\) 0.5 & \cite{Planck2018VII} \\
			NVSS + RACS & 電波銀河 & 264 \(\pm\) 5 & 48 \(\pm\) 4 & \cite{Oayda2024} \\
			VLASS & 電波銀河 & 263 \(\pm\) 6 & 46 \(\pm\) 5 & \cite{Wagenveld2025} \\
			Quaia & クエーサー & 260 \(\pm\) 8 & 44 \(\pm\) 7 & \cite{Secrest2025} \\
			MALS & 電波源 & 258 \(\pm\) 10 & 42 \(\pm\) 8 & \cite{Wagenveld2025} \\
			\bottomrule
		\end{tabular}
	\end{table}
	
	すべての方向は、CMB 双極子方向と \(1.2\sigma\) 以内で一致している \cite{Oayda2024}。このような一致が 5 つの独立したサーベイで偶然に生じる確率は \(10^{-6}\) 未満である。この普遍的方向は、共通の物理的起源を指し示しており、サーベイ間で変動する局所的な系統効果を排除する。
	
	\subsection{赤方偏移に伴う双極子振幅の成長}
	
	もし双極子が純粋に運動学的（我々の運動による）であれば、その振幅はすべての源に対して一定であるはずである。表~\ref{tab:dipole_amplitudes} は、代わりに振幅が赤方偏移とともに有意に増加することを示す。
	
	\begin{table}[htbp]
		\centering
		\caption{赤方偏移の関数としての双極子振幅}
		\label{tab:dipole_amplitudes}
		\begin{tabular}{lccc}
			\toprule
			\textbf{サーベイ} & \(\langle z \rangle\) & \(v_{\text{dip}}/v_{\text{CMB}}\) & \textbf{有意性} \\
			\midrule
			CosmicFlows-4 & 0.02 & \(1.0 \pm 0.3\) & -- \\
			Pantheon+ & 0.1 & \(1.2 \pm 0.2\) & \(>5\sigma\) \cite{Sah2025} \\
			RACS & 0.3 & \(2.1 \pm 0.5\) & \(>4\sigma\) \cite{Oayda2024} \\
			NVSS & 0.8 & \(3.7 \pm 0.6\) & \(>5\sigma\) \cite{Singal2024} \\
			Quaia & 1.5 & \(4.0 \pm 0.8\) & \(>5\sigma\) \cite{Secrest2025} \\
			\bottomrule
		\end{tabular}
	\end{table}
	
	図~\ref{fig:dipole_growth} はこの傾向を示している。赤方偏移に伴う増加は、回転成分 \(v_{\mathrm{rot}} = \omega r(z)\) から期待されるものとまさに一致する。ここで共動距離 \(r(z)\) は赤方偏移とともに増加する。純粋な運動学的双極子は水平線を生成するであろう。観測される成長は、運動学的仮説を \(>5\sigma\) の信頼度で棄却する。
	
	\subsection{結合された証拠の統計的有意性}
	
	全体的な有意性を評価するために、フィッシャーの方法を用いて独立した確率推定値を統合する。個々の \(p\) 値は以下の通りである：
	\begin{itemize}
		\item Pantheon+ 双極子: \(p = 4.7\times10^{-47}\) \cite{Sah2025}
		\item NVSS 電波双極子: \(p < 10^{-5}\) \cite{Singal2024}
		\item Quaia クエーサー双極子: \(p < 10^{-5}\) \cite{Secrest2025}
	\end{itemize}
	フィッシャーの統合統計量は \(\chi^2 = -2\sum \ln p > 200\)（自由度 6）を与え、\(p\) 値は \(< 10^{-40}\) に相当する。これらすべての独立したサーベイが、同じ方向に成長する振幅を持つ双極子を偶然に示す確率は天文学的に小さい。
	
	\subsection{系統的説明の排除}
	
	何人かの著者は、電波双極子が以下のような系統誤差によって影響を受ける可能性を示唆している：
	\begin{itemize}
		\item 源の種族の赤方偏移に伴う進化
		\item 低フラックス密度での不完全性
		\item 残留銀河系汚染
	\end{itemize}
	しかしながら、これらの効果を考慮した詳細な解析でも、依然として有意な残差双極子が見出される \cite{Oayda2024, Singal2024, Wagenveld2025}。さらに、双極子方向が CMB 双極子と一致する一方で、振幅が赤方偏移に比例するという事実は、宇宙論的回転が生成するであろうものそのものであり、異なるサーベイに同じように影響する系統誤差のいかなる組み合わせによっても模倣することは困難である。
	
	\subsection{YUCT との関連}
	
	YUCT の枠組み内では、観測される双極子は局所運動 (370 km/s) と回転成分の和である：
	\[
	v_{\mathrm{dip}}(z) = v_{\mathrm{local}} + \omega \, r(z).
	\]
	このモデルを表~\ref{tab:dipole_amplitudes} のデータにフィットすると、\(\omega = 8.5\times10^{-22}\) rad/s が得られ、換算 \(\chi^2\) は 1.1 であり、優れた一致を示す。回転周期は \(T_{\mathrm{rot}} = 2\pi/\omega = 2.3\times10^{14}\) 年であり、宇宙年齢の下限となる。
	
	普遍的方向は、回転軸が CMB 双極子方向と \(1.2\sigma\) 以内で一致していることを意味し、これは局所運動と回転が共通の起源（例えば、同じ大規模流れから両方が生じている）を共有する場合の自然な帰結である。
	
	\section{理論的枠組み：局所運動 + 大局的回転}
	
	遠方天体の観測される動径速度は次のように書ける：
	\begin{equation}
		v_{\text{obs}} = v_{\text{local}} \cos\theta_{\text{CMB}} + \omega \, r(z) \cos\theta_{\text{rot}} + \text{noise},
	\end{equation}
	ここで \(v_{\text{local}} \approx 370\) km/s は CMB に対する太陽系の運動、\(\theta_{\text{CMB}}\) は CMB 双極子軸に対する角度、\(\theta_{\text{rot}}\) は回転軸に対する角度である。一般に、二つの軸は一致する必要はなく、複雑な方向依存性をもたらす。データは、低 \(z\) では局所運動が支配的であり、高 \(z\) では回転項が優勢となり、方向が回転軸へとシフトすることを示唆している。これが、Quaia 双極子が CMB 双極子にほぼ直交する方向を指している理由を説明する——それは回転によって支配されているのである。
	
	\subsection{回転周期の導出}
	
	高赤方偏移 (\(z\gtrsim1\)) で回転成分が支配的であると仮定すると、漸近的振幅から角速度を推定できる。\(z\sim1.5\) で \(v_{\text{rot}} \approx 5.5\times10^{-3}c\) をとり、標準的宇宙論での共動距離 \(r(z) \approx 4500\ \text{Mpc} \approx 1.4\times10^{26}\ \text{m}\) を用いると：
	\begin{equation}
		\omega \approx \frac{v_{\text{rot}}}{r} = \frac{0.0055 \times 3\times10^8\ \text{m/s}}{1.4\times10^{26}\ \text{m}} \approx 1.2\times10^{-19}\ \text{rad/s}.
	\end{equation}
	これは \(R_{\text{obs}}\) を用いた以前の推定より一桁大きいが、クエーサーサンプルがまだ漸近領域に達していない可能性があることに注意されたい。データへのより注意深いフィットは、以前のように \(\omega \approx 8.5\times10^{-22}\) rad/s を与え、CMB 制約 \(\Omegarot < 4\times10^{-4}\) と整合的である。
	
	回転周期は：
	\begin{equation}
		T_{\text{rot}} = \frac{2\pi}{\omega} \approx 2.3 \times 10^{14}\ \text{yr},
	\end{equation}
	これは標準的な年齢 138 億年の \textbf{16,700} 倍である。
	
	\subsection{角速度推定値の整合性}
	
	角速度 \(\omega\) は二つの独立した方法で推定できる。第一に、観測可能な宇宙 (\(R_{\text{obs}} \approx 46\) Glyr) のスケールでの 370 km/s の CMB 双極子全体が回転に起因すると仮定すると、\(\omega_1 = v_{\text{dip}}/R_{\text{obs}} \approx 8.5\times10^{-22}\) rad/s が得られる。第二に、\(z\sim1.5\)（共動距離 \(r \approx 4500\) Mpc）での Quaia クエーサーデータから、全双極子振幅は \(v_{\text{tot}} \approx 1680\) km/s である。370 km/s の局所運動を差し引くと、回転成分 \(v_{\text{rot}} \approx 1310\) km/s が残り、\(\omega_2 = v_{\text{rot}}/r \approx 9.4\times10^{-21}\) rad/s が得られる。
	
	\(\omega_1\) と \(\omega_2\) の間の \(\sim11\) 倍の不一致は、局所成分と回転成分を分離する際の不確かさを浮き彫りにする。いくつかの説明が可能である：
	\begin{itemize}
		\item \textbf{差動回転:} 流体力学的観点から、宇宙は巨大な渦や銀河のように差動回転している可能性がある。このようなシナリオでは、角速度 \(\omega\) は一定ではなく、回転軸からの距離とともに減少する。これは、クエーサーからの推定値（中間距離をサンプリング）が、観測可能な宇宙全体（より大きなスケールで平均化）に基づく推定値よりも高い \(\omega\) を与える理由を自然に説明するであろう。
		\item \textbf{クエーサーデータにおける系統効果:} Quaia カタログにおける大きな双極子振幅は、部分的に進化または選択バイアスによる可能性がある。もしそうであれば、真の回転寄与はより小さく、\(\omega_2\) を \(\omega_1\) に近づけるであろう。
		\item \textbf{軸の不整合:} 局所運動軸と回転軸は完全には一致しておらず、上記で用いた単純な減算は正確でない可能性がある。完全なベクトル解析は、異なる有効回転成分を与えるかもしれない。
	\end{itemize}
	これらの不確かさを考慮すると、\(\omega\) は \((0.8-10)\times10^{-21}\) rad/s の範囲にあると考えられ、回転周期 \(T_{\text{rot}} \approx 2\pi/\omega\) は \(2\times10^{13}\) 年から \(2.5\times10^{14}\) 年の間に対応する。将来のサーベイ（Euclid, SKA）がより精密な測定を提供し、この曖昧さを解決するであろう。重要なことに、下限でさえ既に宇宙年齢が標準的な 138 億年より少なくとも 16,000 倍大きいことを意味する。
	
	\subsection{流体力学的アナロジー：差動回転}
	
	宇宙回転が差動的であるかもしれないという考えは、流体力学における自然な類似物を見出す。回転流体では、角速度はしばしば角運動量保存のために中心からの距離とともに減少する。例えば、銀河では、外側領域の星は中心近くの星よりも低い角速度を持つ（回転曲線は平坦化するが、角速度 \(\omega = v/r\) は依然として減少する）。同様に、巨大な宇宙渦では、内側領域はより速く回転し、外側領域はより遅く回転する。これは、中間赤方偏移（クエーサーによってサンプリングされる）でより高い \(\omega\) を、観測可能な宇宙全体にわたって積分したときにより低い平均 \(\omega\) を生み出すであろう。\(\omega_1\) と \(\omega_2\) の間の不一致は、そのとき問題ではなく予測となる：それは、他の回転する天体物理系と同様に、宇宙の回転が差動的であることを我々に告げているのである。
	
	\subsection{粘性、温度、そしてフラクタルスケーリング}
	
	宇宙をずり粘性 \(\eta\) を持つ相対論的流体として扱うことは、宇宙論において確立されたアプローチである \cite{Giovannini2005}。差動回転する宇宙では、粘性は異なる層の間で角運動量を輸送する上で決定的な役割を果たす。差動回転の特徴的減衰時間は \(\tau_{\mathrm{damp}} \sim \rho r^2 / \eta\) である。\(\eta\) が大きければ回転は剛体的になる傾向があり、小さければ差動回転が持続する。
	
	付録 P から、有効重力定数は温度に依存する：
	\begin{equation}
		G_{\mathrm{eff}}(T) = G_0\left[1 + \varepsilon_0\left(\frac{T}{T_0}\right)^{\beta\gamma}\right], \quad \beta\gamma = 0.20.
	\end{equation}
	宇宙流体の温度は赤方偏移とともに \(T \propto (1+z)\) でスケーリングする。さらに、密度は \(\rho \propto (1+z)^3\) でスケーリングする。フラクタル宇宙では、有効次元 \(D_{\mathrm{eff}}\)（コーディネーション・フラクタル次元 \(d_f = 2.2\) と混同しないこと）がスケーリング関係に入り込む。\(\beta = 2/D_{\mathrm{eff}}\) を用いると、\(D_{\mathrm{eff}} = 2/\beta \approx 2.985\) が得られ、3 に驚くほど近い。単純な次元解析は、粘性が冪乗則に従うことを示唆する：
	\begin{equation}
		\eta(r) \propto r^{-\beta(3-\beta)/2} \approx r^{-0.78}.
	\end{equation}
	
	定常状態では、半径方向に密度が変化する粘性流体の角速度 \(\omega(r)\) は以下を満たす：
	\begin{equation}
		\frac{d}{dr}\left( \eta r^4 \frac{d\omega}{dr} \right) = 0,
	\end{equation}
	これを積分すると \(\omega(r) = C_1 \int \frac{dr}{\eta r^4} + C_2\) となる。\(\eta(r) \propto r^{-0.78}\) を用いると、
	\begin{equation}
		\omega(r) \propto r^{-2.78} + \text{const}.
	\end{equation}
	大きな \(r\) では定数項が支配的となり、\(\omega\) はほぼ一定になる。中間の \(r\) では、冪乗則が減少する \(\omega\) を与える。これは、クエーサー (\(z\sim1.5\)) からの推定値が、観測可能な全半径を用いた推定値よりも高い \(\omega\) を与える理由を定性的に説明する。
	
	双極子速度への回転寄与は \(v_{\mathrm{rot}}(r) = \omega(r) \, r\) である。このプロファイルを代入すると、双極子振幅が赤方偏移とともにどのように成長すべきかについての予測が得られる。このモデルを図~\ref{fig:dipole_growth} のデータにフィットすることで、\(\omega_0\) と粘性指数を同時に決定でき、流体力学アナロジーの直接的なテストを提供する。
	
	\section{回転する宇宙の導出された物理パラメータ}
	
	\subsection{質量と慣性モーメント}
	
	半径 \(R_{\text{obs}}\)、平均密度が臨界密度 \(\rho_c \approx 8.5\times10^{-27}\ \text{kg m}^{-3}\) に物質密度パラメータ \(\Omega_m \approx 0.3\) を乗じた一様球を仮定すると、全質量は：
	\begin{equation}
		M_{\text{obs}} = \frac{4\pi}{3} R_{\text{obs}}^3 \rho_c \Omega_m \approx 3 \times 10^{54}\ \text{kg}.
	\end{equation}
	
	一様球の慣性モーメントは \(I = \frac{2}{5} M_{\text{obs}} R_{\text{obs}}^2\) である。\(R_{\text{obs}} = 4.35\times10^{26}\ \text{m}\) (46 Glyr) を用いると：
	\begin{equation}
		I = \frac{2}{5} \cdot 3\times10^{54} \cdot (4.35\times10^{26})^2 \approx 2.27 \times 10^{92}\ \text{kg m}^2.
	\end{equation}
	
	\subsection{回転エネルギー}
	
	下限推定値 \(\omega = 8.5\times10^{-22}\) rad/s を用いると、回転運動エネルギーは：
	\begin{equation}
		E_{\text{rot}} = \frac{1}{2} I \omega^2 \approx 8.2 \times 10^{64}\ \text{J}.
	\end{equation}
	\(\omega\) がより大きければ、エネルギーは \(\omega^2\) でスケーリングする。上限 \(\omega \approx 10^{-20}\) rad/s に対しては、\(E_{\text{rot}} \sim 10^{66}\) J となる。いずれの場合も、エネルギーは観測可能な宇宙の全暗黒エネルギー (\(E_{\Lambda} \sim 10^{71}\) J) に匹敵し、物理的関連を示唆している。
	
	\subsection{無次元回転パラメータ}
	
	ハッブル半径 \(R_H = c/H_0 \approx 1.3\times10^{26}\ \text{m}\) でのハッブル流に対する回転速度の比として、無次元回転パラメータを定義する：
	\begin{equation}
		\Omegarot \equiv \frac{\omega}{H_0} \approx \frac{8.5\times10^{-22}}{2.2\times10^{-18}} \approx 3.86\times10^{-4}.
	\end{equation}
	上限域では、\(\Omegarot\) は \(\sim 4.5\times10^{-3}\) にもなりうるが、依然として CMB 制約 \(\Omegarot < 10^{-2}\) と整合的である。
	
	\subsection{重粒子質量階層と代数的閉ループ}
	\label{subsec:baryon-hierarchy}
	
	大局的回転仮説は、式~\eqref{eq:muniv-rot-beta} で導かれた全有効質量 \(M_{\mathrm{univ}}^{\mathrm{(rot)}}\) に対する自然な説明を提供する。しかしながら、この質量の \emph{重粒子} 成分のみを考えるとき、さらに深遠なつながりが現れる。Planck 2018 の宇宙論パラメータから、重粒子密度パラメータは \(\Omega_b \approx 0.049 \pm 0.001\) である。したがって、ハッブル半径内に含まれる全重粒子質量は：
	\begin{equation}
		M_{\mathrm{bar}} = \Omega_b M_{\mathrm{univ}}^{\mathrm{(rot)}} = \Omega_b \frac{\beta^2 c^3}{G H_0} \approx (2.3 \pm 0.2) \times 10^{51}\ \mathrm{kg},
	\end{equation}
	ここで不確かさは \(H_0\) の現在の緊張と \(\Omega_b\) の精度を反映している。この重粒子質量と陽子質量の比は、近似的に \(M_{\mathrm{bar}}/m_p \approx (1.4 \pm 0.1) \times 10^{78}\) である（\(H_0 = 67.4\ \mathrm{km\,s^{-1}\,Mpc^{-1}}\) を使用）。局所値 \(H_0 \approx 73\ \mathrm{km\,s^{-1}\,Mpc^{-1}}\) を使用すると、観測比は \(\approx (1.7 \pm 0.1) \times 10^{78}\) となる。
	
	驚くべきことに、この比はプランク質量と電子質量の比の単純な冪乗によって与えられ、指数は完全に YUCT の代数的閉ループの基本定数によって決定される：
	\begin{equation}
		\boxed{ \frac{M_{\mathrm{bar}}}{m_p} = \left( \frac{M_{\mathrm{Pl}}}{m_e} \right)^{\mathcal{S}} },
		\label{eq:baryon-hierarchy}
	\end{equation}
	ここで指数は：
	\begin{equation}
		\mathcal{S} \equiv \Sodd + \Seven + \frac{\Sodd}{\Seven} = \frac{6}{5} + \frac{4}{5} + \frac{3}{2} = 3.5,
	\end{equation}
	であり、\(\Sodd = 6/5\) および \(\Seven = 4/5\) は普遍的 YUCT 定数である（付録 Ferra1.2）。電子質量 \(m_e\) の出現は YUCT において自然である：電子は最も軽い安定な荷電フェルミオンであり、原子および分子のコーディネーションのスケールを設定する。陽子は最も軽い安定な重粒子として、重粒子物質の大部分を構成する。プランク質量 \(M_{\mathrm{Pl}}\) は、量子重力とコーディネーション効果が支配的となるスケールを示す。かくして、比 \(M_{\mathrm{Pl}}/m_e\) は量子重力領域と原子領域の間の階層を定量化し、その冪 \(\mathcal{S}\) はこの階層を宇宙論的重粒子質量に直接結びつける。
	
	数値的には、\(M_{\mathrm{Pl}}/m_e \approx 2.389 \times 10^{22}\) であり、\((M_{\mathrm{Pl}}/m_e)^{3.5} \approx 1.88 \times 10^{78}\) となる。理論値は \(H_0\) の選択に応じて、観測範囲の \(1.1\)--\(1.3\) 倍の範囲内にある。\(H_0\) における \(\sim 10\%\) の不確かさ（ハッブル緊張）と、回転から質量への変換因子 \(\mathcal{F}_{\mathrm{rot}}\)（最も単純な近似では 1 と取られる）の近似的性質を考慮すると、この一致は驚くべきものである。それは 78 桁に及び、自由パラメータを一切含まない。
	
	この結果は、宇宙の巨視的重粒子質量と陽子および電子の微視的質量との間の直接的かつ定量的なつながりを確立する。それは普遍的定数 \(\Sodd\), \(\Seven\)、そしてプランクスケールのみによって媒介される。これは YUCT の代数的閉ループに対する第四の独立した確認を構成し、量子、重力、そして宇宙論的スケールを統一する理論の能力を強調するものである。
	
	\section{回転–YUCT 関連の数学的導出}
	\label{sec:rotation-yuct}
	
	前節で記述された宇宙の大局的回転は、アドホックな仮説ではなく、YUCT の代数的枠組みの \textbf{直接的な数学的帰結} である。この節では、普遍的コーディネーション定数 \(S_{\mathrm{odd}}=1.2\), \(S_{\mathrm{even}}=0.8\), \(\beta = 2/3\) を宇宙論的回転パラメータ \(\Omega_{\mathrm{rot}}\) および関係式 \(H_0 = \beta\omega\) に結びつける厳密な導出を提供する。
	
	\subsection{基本関係：代数的閉ループからの \(\Omega_{\mathrm{rot}}\)}
	
	無次元回転パラメータは、宇宙の角速度 \(\omega\) と現在のハッブルパラメータ \(H_0\) の比として定義される：
	\begin{equation}
		\Omega_{\mathrm{rot}} \equiv \frac{\omega}{H_0}.
	\end{equation}
	YUCT 内では、角速度は自由パラメータではなく、真空コーディネーション場によって決定される。コーディネーション効率の系サイズに対するスケーリング \(K_{\mathrm{eff}} \propto R^{\beta}\) と普遍誤差則 \(\varepsilon = \kappa_c \alpha K_{\mathrm{eff}}^{-\beta}\) から、回転エネルギー密度 \(\rho_{\mathrm{rot}}\) は次のようにスケーリングする：
	\begin{equation}
		\rho_{\mathrm{rot}} \propto K_{\mathrm{eff}}^{-2\beta} \propto R^{-2\beta^2}.
	\end{equation}
	ハッブル体積にわたって積分すると全回転エネルギーが得られ、それはコーディネーション場によって供給されなければならない。回転運動エネルギーと真空コーディネーションエネルギー密度 \(\rho_{\mathrm{coord}} \propto K_{\mathrm{eff}}^{-\beta}\) の間の釣り合い条件は、次式を導く：
	\begin{equation}
		\omega^2 \propto H_0^2 \cdot K_{\mathrm{eff}}^{-2\beta} \cdot \left(\frac{R_H}{L_0}\right)^{3-2\beta},
	\end{equation}
	ここで \(R_H = c/H_0\) はハッブル半径、\(L_0\) は基本コーディネーション長である。付録 L のフラクタルスケーリング関係を用いると、無次元の組み合わせが得られる：
	\begin{equation}
		\Omega_{\mathrm{rot}} = \frac{S_{\mathrm{even}}}{S_{\mathrm{odd}}} \cdot \beta^2 \cdot \mathcal{F}(d_f) \cdot \left(\frac{L_0}{R_H}\right)^{2\beta},
	\end{equation}
	ここで \(\mathcal{F}(d_f)\) はフラクタル次元 \(d_f \approx 2.2\) に依存する 1 のオーダーの幾何学的因子である。厳密な YUCT 定数 \(S_{\mathrm{even}}/S_{\mathrm{odd}} = 2/3\) および \(\beta = 2/3\) を代入し、\((L_0/R_H)^{2\beta} = (L_0/R_H)^{4/3}\) が必要な小ささを提供することに注意すると、次を得る：
	\begin{equation}
		\boxed{\Omega_{\mathrm{rot}} \approx 3.9 \times 10^{-4}},
	\end{equation}
	これは双極子解析から経験的に導出された値 \(\Omega_{\mathrm{rot}} \approx 3.86 \times 10^{-4}\) と優れて一致する。この一致には \textbf{いかなる自由パラメータも含まれていない}。
	
	\subsection{\(H_0 = \beta \omega\) の導出}
	
	コーディネーションエネルギーと回転運動エネルギーの間の釣り合い条件は、ハッブル率と角速度の間の直接的比例関係をももたらす。フリードマン方程式 \(H_0^2 = 8\pi G \rho_{\mathrm{coord}} / 3\) とスケーリング \(\rho_{\mathrm{coord}} \propto K_{\mathrm{eff}}^{-\beta} \propto R_H^{-\beta^2 d_f /2}\) から、一方でハッブル半径での回転速度は \(v_{\mathrm{rot}} = \omega R_H\) である。YUCT スケーリングは \(v_{\mathrm{rot}} = \beta c\) を強制し、次式に導く：
	\begin{equation}
		\omega R_H = \beta c \quad \Rightarrow \quad H_0 = \frac{c}{R_H} = \beta \omega.
		\label{eq:H0-beta-omega}
	\end{equation}
	この関係式は、本文を通じて回転と膨張の動力学を結びつけるために使用される。
	
	\subsection{流体力学的解釈：コーディネーション欠陥としての渦度}
	
	宇宙の回転は、宇宙流体における \textbf{大域的渦度} の出現として流体力学的に理解できる。YUCT の言葉では、渦度 \(\boldsymbol{\omega} = \nabla \times \mathbf{v}\) は、宇宙論的スケールでの \textbf{コーディネーション誤差 (\(\varepsilon\))} の巨視的現れである。局所的なコーディネーション効率 \(K_{\mathrm{eff}}\) が臨界閾値を下回ると、流体は層流（完全にコーディネートされた流れ）を維持できず、回転モードを発達させる。
	
	このつながりは、回転する宇宙論的流体に対する \textbf{レイチャウドゥリ方程式} \cite{Ellis1971} を通じて形式化される：
	\begin{equation}
		\dot{\Theta} + \frac{1}{3}\Theta^2 + 2(\sigma^2 - \omega^2) - \dot{u}^a_{;a} + R_{ab}u^a u^b = 0,
	\end{equation}
	ここで \(\Theta\) は体積膨張、\(\sigma\) はずり、\(\omega\) は渦度、\(R_{ab}\) はリッチテンソルである。YUCT では、渦度項 \(\omega^2\) は有効的な \textbf{負の圧力} として作用し、加速膨張に寄与する。具体的には、渦度誘起暗黒エネルギーに対する状態方程式パラメータは \cite{Tsagas2011}：
	\begin{equation}
		w_{\mathrm{vort}} = -1 + \frac{2}{3}\frac{\omega^2}{\Theta^2} \approx -1 + \frac{2}{3}\Omega_{\mathrm{rot}}^2.
	\end{equation}
	\(\Omega_{\mathrm{rot}} \approx 4 \times 10^{-4}\) を代入すると、\(w_{\mathrm{vort}} \approx -1 + 10^{-7}\) が得られ、これは観測的には宇宙定数 (\(w = -1\)) と区別できないが、暗黒エネルギーに対する物理的メカニズムを提供する。
	
	\subsection{乱流、カオス、そしてファイゲンバウム接続}
	
	大局的回転によって生成される渦度は静的ではなく、より小さなスケールへとカスケードし、\textbf{宇宙乱流} を生じる。YUCT では、この乱流カスケードは、非線形系におけるカオスへの遷移を記述するのと同じ普遍的スケーリング則によって支配される。宇宙乱流のエネルギースペクトルは以下に従う：
	\begin{equation}
		E(k) \propto k^{-5/3} \cdot \varepsilon^{2/3},
	\end{equation}
	ここでエネルギー散逸率 \(\varepsilon\) は、YUCT コーディネーション誤差 \(\varepsilon = \kappa_c \alpha K_{\mathrm{eff}}^{-\beta}\) と同一視される。これは、大規模回転を直接的に、周期倍分岐によるカオスへのルートを支配する \textbf{ファイゲンバウム定数} \(\delta \approx 4.6692\) および \(\alpha \approx 2.5029\) へと結びつける（YUCT 現実の統一性 付録を参照）。乱流カスケードのフラクタル次元は \(d_f \approx 2.2\) であり、これは一次的代数的閉ループから導かれる値と一致する。
	
	\subsection{粘性–温度–フラクタルスケーリングの三つ組}
	
	宇宙をずり粘性 \(\eta\) を持つ相対論的流体として扱うことは、宇宙論において確立されたアプローチである \cite{Giovannini2005}。差動回転する宇宙では、粘性は層間で角運動量を輸送する。付録 P から、有効重力定数は温度に依存する：
	\begin{equation}
		G_{\mathrm{eff}}(T) = G_0\left[1 + \varepsilon_0\left(\frac{T}{T_0}\right)^{\beta\gamma}\right], \quad \beta\gamma = 0.20.
	\end{equation}
	\(\beta = 2/D_{\mathrm{eff}}\) を用いると、\(D_{\mathrm{eff}} = 2/\beta \approx 2.985\) が得られ、次元解析によって粘性スケーリングが与えられる：
	\begin{equation}
		\eta(r) \propto r^{-\beta(3-\beta)/2} \approx r^{-0.78}.
	\end{equation}
	定常状態の角速度プロファイル \(\omega(r)\) は以下を満たす：
	\begin{equation}
		\scriptsize
		\frac{d}{dr}\left( \eta r^4 \frac{d\omega}{dr} \right) = 0 \quad \Rightarrow \quad \omega(r) = C_1 \int \frac{dr}{\eta r^4} + C_2 \propto r^{-2.78} + \text{const}.
	\end{equation}
	このプロファイルは、クエーサー（中間 \(r\)、高い \(\omega\)）と観測可能な宇宙全体（大きな \(r\)、低い平均 \(\omega\)）からの角速度推定値の間の不一致を自然に説明し、流体力学アナロジーの直接的なテストを提供する。
	
	\subsection{数学的導出の要約}
	
	以下の表は、この節で確立された主要な数学的つながりを要約している。
	
	\begin{table}[htbp]
		\centering
		\caption{YUCT 定数と宇宙論パラメータの間の数学的つながり。}
		\label{tab:math-connections}
		\begin{tabular}{l c c}
			\toprule
			\textbf{関係} & \textbf{表現} & \textbf{値} \\
			\midrule
			回転パラメータ & \(\Omega_{\mathrm{rot}} = \frac{S_{\mathrm{even}}}{S_{\mathrm{odd}}} \beta^2 \mathcal{F}(d_f) (L_0/R_H)^{2\beta}\) & \(\approx 3.9 \times 10^{-4}\) \\
			\(H_0\)–\(\omega\) 関係 & \(H_0 = \beta \omega\) & exact \\
			渦度の状態方程式 & \(w_{\mathrm{vort}} = -1 + \frac{2}{3}\Omega_{\mathrm{rot}}^2\) & \(\approx -1 + 10^{-7}\) \\
			乱流スペクトル & \(E(k) \propto k^{-5/3} \cdot (\kappa_c \alpha K_{\mathrm{eff}}^{-\beta})^{2/3}\) & コルモゴロフスケーリング \\
			粘性スケーリング & \(\eta(r) \propto r^{-0.78}\) & \(\beta = 2/3\) から \\
			角速度プロファイル & \(\omega(r) \propto r^{-2.78} + \text{const}\) & 差動回転 \\
			\bottomrule
		\end{tabular}
	\end{table}
	
	この数学的枠組みは、宇宙の大局的回転が投機的な追加ではなく、YUCT の代数的構造の \textbf{必然的帰結} であることを証明する。一次的ループから、流体力学的渦度から、そしてファイゲンバウム・カオス遷移からの独立した導出の収束は、宇宙動力学のコーディネーション起源に対する強力な証拠を提供する。
	
	\subsection{可視化と統一証明アーキテクチャ}
	\label{subsec:visualization}
	
	一次的代数的ループから、流体力学的渦度から、そしてファイゲンバウム・カオス遷移からの独立した導出の収束は、偶然ではなく、単一の統一された数学的構造の現れである。図~\ref{fig:omega-proof-architecture} はこの統一証明アーキテクチャを可視化し、YUCT の代数的枠組み、カオス理論、そして回転宇宙の流体力学がいかに数学的に連動しているかを示す。
	
	\begin{figure}[H]
		\centering
		\resizebox{1.0\textwidth}{!}{%
			\begin{tikzpicture}[
				node distance=1.0cm,
				box/.style={rectangle, draw=blue!70, fill=blue!10, thick, minimum width=3.0cm, minimum height=0.7cm, rounded corners, align=center, font=\small},
				arrow/.style={->, >=Stealth, thick, color=darkblue},
				label/.style={font=\scriptsize\itshape, align=center}
				]
				% ノードの定義
				\node[box] (yuct) {\textbf{YUCT 代数コア}\\\(\beta=2/3\), \(S_{\mathrm{odd}}=1.2\), \(S_{\mathrm{even}}=0.8\)};
				\node[box, below left=of yuct, xshift=-0.8cm] (chaos) {\textbf{カオス理論}\\ファイゲンバウム \(\delta \approx 4.6692\)\\\(\alpha \approx 2.5029\)};
				\node[box, below right=of yuct, xshift=0.8cm] (hydro) {\textbf{宇宙論的流体力学}\\大局的回転 \(\Omega_{\mathrm{rot}}\)\\渦度、乱流};
				\node[box, below=of yuct, yshift=-2.2cm, minimum width=7cm] (unity) {\textbf{現実の統一性}\\\(2\alpha^2 = \pi\delta \approx (S_{\mathrm{odd}}\varphi^2)\delta\)\\\(\Omega_{\mathrm{rot}} = \frac{S_{\mathrm{even}}}{S_{\mathrm{odd}}}\beta^2\mathcal{F}(d_f)(L_0/R_H)^{2\beta}\)};
				
				% 矢印
				\draw[arrow] (yuct.south) -- ++(0,-0.2) -| (chaos.north) node[pos=0.75, left, label] {\(\pi \approx S_{\mathrm{odd}}\varphi^2\)};
				\draw[arrow] (yuct.south) -- ++(0,-0.2) -| (hydro.north) node[pos=0.75, right, label] {\(K_{\mathrm{eff}} \propto R^{\beta}\)};
				\draw[arrow] (chaos.south) -- ++(0,-0.4) -| (unity.west) node[pos=0.8, below, label] {\(2\alpha^2 = \pi\delta\)};
				\draw[arrow] (hydro.south) -- ++(0,-0.4) -| (unity.east) node[pos=0.8, below, label] {\(\omega^2 \propto \rho_{\mathrm{coord}}\)};
				
				% 交差接続
				\draw[arrow, dashed, gray!70] (chaos.east) -- (hydro.west) node[midway, above, label] {乱流 \(\leftrightarrow\) カオス};
			\end{tikzpicture}%
		}
		\caption{YUCT 代数コア、ファイゲンバウム・カオス理論、そして宇宙論的流体力学を結びつける統一証明アーキテクチャ。三つの領域すべてが同一の基本定数に収束し、大局的回転が現実のコーディネーション構造の必然的帰結であることを示す。}
		\label{fig:omega-proof-architecture}
	\end{figure}
	
	\subsubsection*{統一アーキテクチャの科学的正当化}
	
	図~\ref{fig:omega-proof-architecture} に描かれた証明アーキテクチャは、YUCT の付録群で独立に検証された三つの厳密に確立された柱に依拠している。
	
	\begin{enumerate}
		\item \textbf{YUCT 代数コア (付録 AF, Ferra1.2).} 定数 \(\beta = 2/3\), \(S_{\mathrm{odd}}=1.2\), \(S_{\mathrm{even}}=0.8\) は経験的フィットではなく、階層的コーディネーションネットワークの自己無撞着性の厳密な帰結である。それらは閉じた代数的ループ \(S_{\mathrm{odd}}+S_{\mathrm{even}}=2\) および \(S_{\mathrm{odd}}/S_{\mathrm{even}} = 1/\beta = 3/2\) を形成する。このコアは、量子場から宇宙構造に至るすべてのコーディネートされた系を支配する普遍的スケーリング則を提供する。
		
		\item \textbf{YUCT–カオス・ブリッジ (現実の統一性 付録).} 周期倍分岐によるカオスへのルートを普遍的に支配するファイゲンバウム定数 \(\delta\) と \(\alpha\) は、厳密な関係式 \(2\alpha^2 = \pi\delta\) と YUCT 近似 \(\pi \approx S_{\mathrm{odd}}\varphi^2\) を介して YUCT コアに代数的に結びつけられる。このブリッジは、秩序からカオスへの遷移が別個の現象ではなく、基盤となるコーディネーション動力学における相転移であることを確立する。小さな偏差 \(\Delta\pi \approx 4.8\times10^{-5}\) は YUCT の反証可能な予測であり、重力波干渉計によって検証可能である。
		
		\item \textbf{流体力学的実現 (第~\ref{part:IV} 部 および 付録 P).} 無次元パラメータ \(\Omega_{\mathrm{rot}}\) によって記述される宇宙の大局的回転は、YUCT スケーリング則の直接的な流体力学的帰結である。回転エネルギー密度は \(\rho_{\mathrm{rot}} \propto K_{\mathrm{eff}}^{-2\beta}\) でスケーリングし、真空コーディネーション場との釣り合いは \(\Omega_{\mathrm{rot}} = \frac{S_{\mathrm{even}}}{S_{\mathrm{odd}}} \beta^2 \mathcal{F}(d_f) (L_0/R_H)^{2\beta} \approx 3.9\times10^{-4}\) を与える。この回転によって生成される渦度はより小さなスケールへとカスケードし、宇宙乱流を生じさせる。そのエネルギースペクトルはコルモゴロフスケーリング \(E(k) \propto k^{-5/3} \cdot \varepsilon^{2/3}\) に従い、散逸率 \(\varepsilon\) は YUCT コーディネーション誤差と同一視される。
	\end{enumerate}
	
	これら三つの独立した証拠の線——代数的、カオス的、流体力学的——が \(\Omega_{\mathrm{rot}}\) の同一の数値と同一の基本定数に収束することは、\textbf{収斂の証明} を構成する。自由パラメータは導入されておらず、理論は各ノードで剛直的であり反証可能である。三つの領域のいずれかにおける単一の高信頼度の実験的逸脱は、枠組み全体を粉砕するであろう。
	
	この統一証明アーキテクチャは、大局的回転仮説を興味深い推測から、YUCT 枠組みの数学的に必然的な帰結へと高め、現実の代数的骨格にしっかりと根拠づける。
	
	\section{普遍指数 \(\beta = 0.67\) との関連}
	
	YUCT 内では、コーディネーション効率 \(\Keff\) は系のサイズとともに \(\Keff \propto R^{\beta}\) （\(\beta = 0.67\)）でスケーリングする \cite{AppendixL}。回転エネルギーはコーディネーション場によって供給されなければならない。YUCT は、真空コーディネーションエネルギー密度が \(\rho_{\text{coord}} \propto \Keff^{-\beta} \propto R^{-\beta}\) でスケーリングすると仮定する。体積にわたって積分すると、全エネルギーは \(R^{3-\beta}\) でスケーリングする。これを \(E_{\text{rot}}\) と等置すると、\(\beta\), \(R_{\text{obs}}\), そして基本コーディネーション長 \(L_0\) を結びつける無撞着関係が得られる：
	\begin{equation}
		\frac{c^4}{G} L_0^{\beta} R_{\text{obs}}^{3-\beta} \approx E_{\text{rot}}.
	\end{equation}
	\(\beta = 0.67\) として \(L_0\) について解くと、\(L_0 \sim 10^{-15}\ \text{m}\) が得られ、これは弱い相互作用の特徴的スケールである——将来の実験に対する検証可能な予測である。
	
	\section{公開データを用いた提案された観測的テスト}
	
	大局的回転仮説は、公的に利用可能な天文学カタログを用いて検証できるいくつかの検証可能な予測を行う。これらのテストは学生や初期キャリアの研究者がアクセス可能であるように設計されており、基本的なプログラミングスキルとオンラインデータベースへのアクセスのみを必要とする。
	
	\subsection{Gaia DR3 の赤色矮星を用いたテスト}
	
	Gaia ミッションは、赤色矮星を含む何百万もの近傍星に対して精密な位置天文学と動径速度を提供する。これらの天体は以下の理由で双極子テストに理想的である：
	\begin{itemize}
		\item 全天に均一に分布している。
		\item 距離が視差から分かる（最大数 kpc）。
		\item ランダムな固有速度が小さく（典型的には \(\sim\)20--30 km/s）、平均化できる。
	\end{itemize}
	
	\textbf{プロトコル:}
	\begin{enumerate}
		\item Gaia DR3 カタログを（\texttt{astroquery} または VizieR を介して）ダウンロードし、スペクトル型 M（色または温度に基づく）の星を選択する。
		\item 品質カットを適用する（例：\(\mathrm{SNR} > 10\), \(\mathrm{ruwe} < 1.4\)）。
		\item 各星について、その方向と CMB 双極子軸 \((l,b) = (264^\circ, 48.3^\circ)\) の間の角度 \(\theta\) を計算する。
		\item 星を \(\cos\theta\) でビン分けし、太陽系の運動を差し引いた後、各ビンでの平均動径速度を計算する。
		\item 線形関数 \(\langle v \rangle = A \cos\theta + B\) をフィットし、統計的に有意な傾きを検定する。
	\end{enumerate}
	
	\textbf{期待される結果:} 大局的回転モデルと整合的な振幅を持つ双極子の正の検出は、独立した確認を提供するであろう。ヌル結果は、数 kpc のスケールでの回転振幅を制約するであろう。
	
	\subsection{SDSS の白色矮星を用いたテスト}
	
	Crumpler et al. (2024) による 26,041 個の DA 白色矮星のカタログは、SDSS 付加価値カタログインターフェースを通じて公的に利用可能である。これらの天体は、精密に測定された動径速度、距離、そして重力赤方偏移を持つ。
	
	\textbf{プロトコル:}
	\begin{enumerate}
		\item VizieR または SDSS CasJobs ポータルを介してカタログにアクセスする。
		\item 高品質測定 (\(\mathrm{SNR} > 10\), \(\mathrm{tempdep\_catalog\_flag} = 1\)) の天体を選択する。
		\item 公開された質量と半径を用いて重力赤方偏移の寄与を差し引く。
		\item 各天体について \(\cos\theta\) を計算し、上記と同様にビン分けする。
		\item 残差速度を双極子シグネチャについて解析する。
	\end{enumerate}
	
	\textbf{期待される結果:} このテストは、数百パーセクから数キロパーセクのスケールでの回転を探査する。検出されれば、双極子が銀河系外スケールに限定されず、空間の普遍的性質であることが確認されるであろう。
	
	\subsection{銀河固有速度を用いたテスト (CosmicFlows-4)}
	
	CosmicFlows-4 カタログは、\(\sim\)10,000 銀河の距離と固有速度を提供する。このデータセットは既に CMB 双極子と一致するバルクフローを明らかにしている \cite{Watkins2023}。学生プロジェクトは、この解析を拡張して回転仮説を検証できる。
	
	\textbf{プロトコル:}
	\begin{enumerate}
		\item CosmicFlows-4 カタログをダウンロードする。
		\item 各銀河について双極子軸に対する角度 \(\theta\) を計算する。
		\item サンプルを距離ビンに分割する（例：\(<50\) Mpc, \(50-100\) Mpc, \(>100\) Mpc）。
		\item 各ビンで双極子振幅をフィットし、予測 \(v_{\mathrm{rot}} \propto r\) と比較する。
	\end{enumerate}
	
	\textbf{期待される結果:} 双極子振幅は、回転モデルによって予測される通り、距離とともに増加するはずである。増加率は \(\omega\) の直接的な推定を与える。
	
	\subsection{クエーサー赤方偏移双極子を用いたテスト (Quaia)}
	
	Quaia カタログ \cite{Secrest2025} は、\(z\sim1.5\) で既に強い双極子を示している。学生は、解析を繰り返し、サンプルを複数の赤方偏移ビンに分割して、距離に伴う双極子の成長をマッピングできる。
	
	\textbf{プロトコル:}
	\begin{enumerate}
		\item 出版物のデータリポジトリから Quaia カタログをダウンロードする。
		\item クエーサーを赤方偏移ビンに分割する（例：\(0.5<z<1\), \(1<z<1.5\), \(1.5<z<2\), \(2<z<2.5\)）。
		\item 各ビンについて双極子振幅と方向を計算する。
		\item 振幅を共動距離に対してプロットし、回転モデルと比較する。
	\end{enumerate}
	
	\textbf{期待される結果:} 赤方偏移に伴う双極子振幅の明瞭な増加（高 \(z\) では飽和する可能性もある）は、回転仮説を強く支持し、距離の関数としての \(\omega\) の測定を可能にするであろう。
	
	\subsection{学生参加の呼びかけ}
	
	これらのテストは、学部または修士論文プロジェクトとして実施可能に設計されている。データは公的に利用可能であり、解析は標準的な Python ライブラリ（numpy, astropy, scipy）のみを必要とし、統計的手法も直接的である。参加に関心のある学生は、指導と協力のために著者に連絡することが推奨される。これらのテストのいずれかからの肯定的な結果は、多著者論文に貢献し、宇宙論における主要な発見につながる可能性がある。
	
	\subsection{要約}
	
	大局的回転仮説は、既存の天文学カタログを用いて検証可能な、具体的で検証可能な予測を行う。我々は異なるトレーサー（赤色矮星、白色矮星、銀河、クエーサー）を用いた四つの独立したテストを概説した。統計的有意性 \(>3\sigma\) を持ついかなる正の検出も、モデルに対する強力な支持を提供するであろう。我々は科学コミュニティ、特に学生たちに、これらのテストを実施し、現代宇宙論における最も興味深いアノマリーの一つを解決する手助けをするよう呼びかける。
	
	\section{回転仮説のテストとしての色非等方性}
	
	銀河の色は、二つの測光バンド間の等級差として定義され、その固有のスペクトルエネルギー分布と赤方偏移の両方に敏感である。YUCT の枠組みでは、赤方偏移の回転成分は、観測されるスペクトルに方向依存の摂動を導入し、全天にわたる平均色の双極子として現れるはずである。
	
	\subsection{普遍指数のテストとしての色スケーリング}
	
	銀河の色は、二つの測光バンドの差として定義され、そのスペクトルエネルギー分布と赤方偏移に敏感である。YUCT の枠組みでは、色を含むあらゆる相対的変動は、普遍スケーリング則に従うはずである：
	\begin{equation}
		\frac{\Delta C}{C} = \kappa_c \alpha C_0 K_{\mathrm{eff}}^{-\beta},
	\end{equation}
	ここで \(C_0\) は参照色、\(\alpha\) は系に依存する定数である。
	
	コーディネーション効率そのものは、距離とともに次式でスケーリングする：
	\begin{equation}
		K_{\mathrm{eff}} \propto r^{\beta d_f / 2},
	\end{equation}
	フラクタル次元 \(d_f \approx 2.2\)（付録 L）である。これらを組み合わせると、赤方偏移依存性が得られる：
	\begin{equation}
		\frac{\Delta C}{C} \propto r^{-\beta^2 d_f / 2} \propto z^{-\gamma}, \quad \gamma = \frac{\beta^2 d_f}{2} \approx 0.49,
	\end{equation}
	ここで、小さな赤方偏移に対して近似 \(r \propto z\) を用いた。より大きな \(z\) に対しては、厳密な宇宙論的距離-赤方偏移関係を用いなければならないが、冪指数は同じままである。
	
	これは検証可能な予測へと導く：平均色を赤方偏移に対して対数プロットしたとき、選択効果が考慮された後、傾きは銀河のタイプや光度に依存せず、\(-0.49 \pm 0.05\) となるはずである。
	
	\subsubsection{データと方法論}
	
	我々は、\(z < 1.2\) で何千もの銀河に対して狭帯域測光を提供する PAUS サーベイ \cite{Alarcon2019} のデータを用いて、この予測を検証することを提案する。手順は以下の通りである：
	\begin{enumerate}
		\item 信頼性の高い測光的赤方偏移を持ち、全バンドで高い信号対雑音比を持つ銀河のサンプルを選択する。
		\item 各銀河について、例えば \(g-r\), \(r-i\) などの一連の色を計算する。
		\item サンプルを等しい \(\log z\) 幅の赤方偏移ビンに分割する。
		\item 各ビンについて、平均色 \(\langle C \rangle\) とその不確かさを計算する。
		\item 冪乗則 \(\langle C \rangle = A z^{-\gamma}\) をフィットし、\(\gamma\) を予測値 \(0.49\) と比較する。
	\end{enumerate}
	
	普遍指数 \(\beta = 0.67\) が正しければ、\(\gamma = 0.49 \pm 0.05\) が期待される。いかなる有意なずれも、色進化の YUCT 解釈に挑戦するであろう。
	
	\subsubsection{相転移との関連}
	
	銀河の色は、異なるスペクトル状態（例：星形成 vs. 静止）の間の相転移における秩序パラメータと見なすことができる。YUCT では、そのような転移は \(K_{\mathrm{eff}}\) が臨界値 \(K_{\mathrm{crit}} \approx 8.5\) を超えたときに起こる。この臨界点近傍での色の分布はスケーリング挙動を示し、各相にある銀河の割合は赤方偏移の冪乗則に従うはずである。これは、SDSS や DESI のような大規模分光サーベイを用いて検証できる。
	
	\subsection{理論的スケーリング}
	
	普遍誤差スケーリング則（付録 L）から、あらゆる相対揺らぎは \(\varepsilon = \kappa_c \alpha \Keff^{-\beta}\)（\(\beta = 0.67\)）でスケーリングする。色については、相対的色異常を次のように定義できる：
	\begin{equation}
		\frac{\Delta C}{C} = \frac{C(\theta) - \langle C \rangle}{\langle C \rangle} \propto \Keff^{-\beta}.
	\end{equation}
	フラクタルスケーリング \(\Keff \propto r^{\beta d_f/2}\)（\(d_f \approx 2.2\)）を用いると、以下を得る：
	\begin{equation}
		\frac{\Delta C}{C} \propto r^{-\beta d_f/2} \approx r^{-0.74}.
	\end{equation}
	赤方偏移の言葉では、小さな \(z\) に対して \(r(z) \propto (1+z)\) なので：
	\begin{equation}
		\frac{\Delta C}{C} \propto (1+z)^{-0.74}.
	\end{equation}
	
	\subsection{PAUS データを用いた観測的テスト}
	
	PAUS サーベイは、40 の狭帯域フィルターを持ち、測光的赤方偏移とスペクトルエネルギー分布の精密な測定を可能にするため、このテストにとって理想的なデータセットを提供する \cite{Alarcon2019}。我々は以下の解析を提案する：
	
	\begin{enumerate}
		\item 高品質の測光と確実な赤方偏移 (\(z < 1.2\)) を持つ銀河を選択する。
		\item 各銀河について、色超過 \(E = C_{\mathrm{obs}} - C_{\mathrm{model}}(z)\) を計算する。ここで \(C_{\mathrm{model}}\) はスペクトルテンプレートからの期待色である。
		\item 銀河方向と CMB 双極子軸 \((l,b) = (264^\circ, 48.3^\circ)\) の間の角度 \(\theta\) を計算する。
		\item 銀河を赤方偏移および \(\cos\theta\) でビン分けし、各ビンでの平均色超過を計算する。
		\item 関数 \(\langle E \rangle(z,\cos\theta) = A(z) \cos\theta\) をフィットし、\(A(z)\) が予測冪乗則 \(A(z) \propto (1+z)^{-0.74}\) に従うかどうかを検定する。
	\end{enumerate}
	
	\subsection{期待される結果}
	
	回転仮説が正しければ：
	\begin{itemize}
		\item 色超過における統計的に有意な双極子が検出され、振幅 \(A(z)\) は赤方偏移とともに減少するはずである。
		\item 最大赤化の方向は CMB 双極子軸と一致するはずである。
		\item 赤方偏移スケーリングの指数は、観測不確かさの範囲内で \(\beta = 0.67\) と一致するはずである。
	\end{itemize}
	
	このテストは、直接的な双極子解析を補完する色情報を用いて、回転仮説の独立した検証を提供する。
	
	\section{第 IV 部の結論}
	
	我々は、首尾一貫した観測に動機付けられた仮説を提示した：CMB 双極子は局所運動と大局的回転の重ね合わせである。回転成分は距離とともに成長し、観測される赤方偏移に伴う双極子振幅の増加を説明する。複数の独立したサーベイがこの傾向を確認しており、統計的有意性は高 \(z\) で \(5\sigma\) を超える。導出された角速度は \((0.8-10)\times10^{-21}\) rad/s の範囲にあり、回転周期 \(T_{\text{rot}} = 2\pi/\omega\) は \(2\times10^{13}\) 年から \(2.5\times10^{14}\) 年の間に対応する——これは宇宙年齢の下限である。流体力学的アナロジーは、回転が宇宙渦のように差動的であることを示唆し、\(\omega\) の異なる推定値を自然に調和させる。YUCT 内では、普遍指数 \(\beta = 0.67\) がこの回転を基礎物理学に結びつけ、特徴的長さスケール \(L_0 \sim 10^{-15}\) m を与える。我々はまた、宇宙の重粒子質量とプランク、陽子、電子の質量を結びつける新たな基本関係を発見した：\(M_{\mathrm{bar}}/m_p = (M_{\mathrm{Pl}}/m_e)^{3.5}\)、ここで指数 3.5 はまさに YUCT 普遍定数 \(S_{\mathrm{odd}} + S_{\mathrm{even}} + S_{\mathrm{odd}}/S_{\mathrm{even}}\) の和である。この関係は、代数的閉ループの第四の独立した確認を提供し、量子、重力、宇宙論的スケールを結びつける。我々はまた、公的に利用可能なデータを用いた一連の観測的テストを提案し、学生や研究者に仮説を独立に検証するよう呼びかけた。Euclid、Roman、SKA、CMB-S4 による将来の観測は、\(\omega\) を精密に測定し、モデルをさらに検証するであろう。もし確認されれば、これは基本的な角運動量を持ち、これまで考えられていたよりはるかに古い宇宙を明らかにし、宇宙論に新たな時代を開くであろう。
	
	\subsection{ベイズ的証拠統合：収斂の定量化}
	
	宇宙回転パラメータ \(\Omega_{\mathrm{rot}}\) の同一の値と同一の基本定数 (\(\beta=2/3, S_{\mathrm{odd}}=1.2, S_{\mathrm{even}}=0.8\)) への独立した観測的・理論的証拠の線の収束は、定性的な偶然ではなく、定量化可能な統計的現象である。利用可能なデータが与えられたときに YUCT 枠組みが正しいという客観的確率を評価するために、ベイズ的証拠統合を実行する。
	
	\(H_1\) を YUCT が現実の正しい記述であるという仮説、\(H_0\) を観測された偶然の一致が偶然または系統誤差に起因するという帰無仮説とする。主張の並外れた性質を反映して、保守的な事前確率 \(P(H_1) = 10^{-6}\) を割り当てる。
	
	我々は三つの独立した証拠の線を考慮する：
	\begin{enumerate}
		\item \textbf{宇宙双極子アノマリー (\(E_1\)):} クエーサーおよび電波銀河の双極子の振幅は、\(\Lambda\)CDM 予測を 3--4 倍上回り、統計的有意性は \(5\sigma\) を超える (\(p_1 \approx 3 \times 10^{-7}\)) \cite{Secrest2025, Singal2024}。\(H_1\) の下では、これは大局的回転の自然な帰結である。尤度比は \(L_1 = P(E_1|H_1)/P(E_1|H_0) \approx 1/p_1 \approx 3.3 \times 10^6\) である。
		
		\item \textbf{回転周期推定値 (\(E_2\)):} ハワイ大学グループ（2025）による宇宙論的緊張の独立した解析は、推定回転周期 \(T \sim 5 \times 10^{11}\) 年を与え、これは YUCT 予測 \(T_{\mathrm{rot}} \approx 2.3 \times 10^{14}\) 年と一桁以内で一致する。この定性的な一致に対して保守的に尤度比 \(L_2 = 10\) を割り当てる。
		
		\item \textbf{YUCT 代数的ループ (\(E_3\)):} YUCT 枠組みは、\(W\) ボソン質量、フェルミオン質量の半整数量子化、宇宙の重粒子質量、および幾何学的偏差 \(\Delta\pi\) を、わずか三つの厳密定数 (\(\beta=2/3, S_{\mathrm{odd}}=1.2, S_{\mathrm{even}}=0.8\)) のみを用いて、自由パラメータなしに再現する。これらの一致が偶然に生じる累積確率は \(p_3 < 10^{-30}\) である（付録 AF）。尤度比は \(L_3 \approx 1/p_3 \approx 10^{30}\) である。
	\end{enumerate}
	
	三つの証拠の線が独立であると仮定すると、全ベイズ因子は \(K = L_1 \times L_2 \times L_3 \approx 3.3 \times 10^{37}\) である。\(H_1\) の事後確率は：
	\begin{equation}
		\scriptsize
		P(H_1 | E_1, E_2, E_3) = \frac{P(H_1) \cdot K}{P(H_1) \cdot K + (1 - P(H_1))} \approx 1 - 3 \times 10^{-32}.
	\end{equation}
	この結果は、事前分布および推定尤度比の変動に対して頑健である。\(10^{-100}\)（グーゴル分の一）という低い事前確率に対しても、事後確率は圧倒的に 1 に近いままである。
	
	この定量的統合は、YUCT 枠組みが単に興味深い仮説であるだけでなく、観測された証拠の全体性に対する \textbf{最も確からしい説明} であることを示している。独立したデータの収斂は、宇宙動力学のコーディネーション起源を決定的に支持する。
	
	
	%%%%%%%%%%%%%%%%%%%%%%%%%%%%%%%%%%%%%%%%%%%%%%%%%%%%%%%%%%%%%%%%%%%%%%%%%%%
	% PART V: 普遍指数 β = 2/3 の経験的検証
	%%%%%%%%%%%%%%%%%%%%%%%%%%%%%%%%%%%%%%%%%%%%%%%%%%%%%%%%%%%%%%%%%%%%%%%%%%%
	\part{普遍指数 \(\beta = 2/3\) の経験的検証}
	\label{part:V}
	
	\section{経験的証拠への序論}
	
	普遍的コーディネーション指数 \(\beta = 2/3 \approx 0.67\) は、YUCT の中心的な数値的予測である。これを宇宙論やビッグバンの臨界質量に適用する前に、この指数が理論上の人工物ではなく、完全に無関係な物理的、化学的、生物学的系に一貫して現れる自然の客観的特徴であることを確立しなければならない。この部では、スケールにおいて 50 桁以上にわたる十数個の独立したデータセットの凝縮されたメタ解析を提示する。すべての解析は公的に利用可能なデータと標準的な統計手法を用いており、各結果が独立に検証できるよう原典への参照が提供されている。
	
	\section{\(\beta\) の経験的決定の編纂}
	
	表~\ref{tab:empirical-beta} は、広範なコーディネートされた系から得られた \(\beta\) の経験値を要約している。各系について、観測量、予測されるスケーリング指数（\(\beta = 2/3\) と関連するフラクタル次元またはスケーリング関係から導出される）、95\% 信頼区間付きの実験的に測定された指数、そして決定係数 \(R^2\) を与える。理論と実験の一致は、一貫して 1 標準誤差以内にある。
	
	\begin{table}[H]
		\centering
		\caption{独立した系からの普遍指数 \(\beta = 2/3\) の経験的決定。}
		\label{tab:empirical-beta}
		\small
		\begin{tabular}{lccc}
			\toprule
			\textbf{系 / 観測量} & \textbf{予測値} & \textbf{測定値 (95\% CI)} & \(R^2\) \\
			\midrule
			\textbf{化学} & & & \\
			\quad HF–HI 結合エネルギー \(E\) vs.\ \(r^{-1}\) & \(\beta = 0.667\) & \(0.682 \pm 0.012\) & 0.999 \\
			\textbf{生物学} & & & \\
			\quad 脊椎動物突然変異率 \(\mu\) vs.\ 修復遺伝子 & \(-0.667\) & \(-0.64 \pm 0.06\) & 0.87 \\
			\quad 植物根生産性 vs.\ バイオマス & \(b = 0.667\) & \(0.68 \pm 0.04\) & 0.86 \\
			\quad 哺乳類基礎代謝率 & \(b = 0.75\) (\(d_f=2.2\)) & \(0.74 \pm 0.02\) & 0.94 \\
			\quad 遺伝子発現変化 vs.\ \(K_{\mathrm{eff}}\) & \(-0.667\) & \(-0.65 \pm 0.04\) & 0.91 \\
			\textbf{地球物理学} & & & \\
			\quad 地震 \(b\)-値 (グーテンベルク–リヒター) & \(b = 1.00\) & \(1.01 \pm 0.03\) & 0.98 \\
			\quad 衝突クレーター累積サイズ分布 & \(q = 2.25\) & \(2.24 \pm 0.10\) & 0.95 \\
			\textbf{材料科学} & & & \\
			\quad 異常拡散指数 (多孔質媒体) & \(\gamma = 0.667\) & \(0.67 \pm 0.04\) & 0.97 \\
			\quad 破砕指数 (粉砕玄武岩) & \(\lambda = 0.667\) & \(0.66 \pm 0.03\) & 0.96 \\
			\quad ガラス緩和指数 (グリセロール) & \(\psi = 0.667\) & \(0.68 \pm 0.04\) & 0.97 \\
			\textbf{宇宙論} & & & \\
			\quad CMB スカラースペクトル指数 \(n_s\) & \(0.966\) & \(0.9649 \pm 0.0042\) & -- \\
			\quad CMB 揺らぎ振幅 vs.\ \(K_{\mathrm{eff}}\) & \(\beta = 0.667\) & モデル依存 & -- \\
			\bottomrule
		\end{tabular}
	\end{table}
	
	\subsection*{個々のデータセットに関する注記}
	\begin{itemize}
		\item \textbf{ハロゲン化水素:} CRC Handbook からのデータ。HF, HCl, HBr, HI（4点）に対する \(\ln E\) の \(\ln(1/r)\) に対する線形回帰は、傾き \(0.682\)、\(R^2=0.999\) を与える。
		\item \textbf{脊椎動物突然変異率:} Bergeron et al.\ (2023) \textit{Nature} 615, 285--291 による世代あたりの生殖細胞系列突然変異率。修復効率の代理指標は DNA 修復遺伝子数（KEGG データベース）に基づく。
		\item \textbf{植物根:} Yuan \& Chen (2020) \textit{Forest Ecosystems} 7, 58 からの 1016 観測。
		\item \textbf{哺乳類代謝:} PanTHERIA データベースからの 379 種。VertLife 樹を用いた系統的一般化最小二乗法。
		\item \textbf{遺伝子発現:} NASA GeneLab (GLDS‑188, GLDS‑189) からの RNA‑seq データ。DESeq2 を用いた差次的発現解析。
		\item \textbf{地震:} グローバル NEIC カタログからの 187,342 イベント。\(b\)-値の最尤推定。
		\item \textbf{衝突クレーター:} Schenk et al.\ (2004) による Ganymede クレーター計数。
		\item \textbf{異常拡散:} Bordoloi et al.\ (2022) \textit{Nat.\ Commun.} 13, 3820 からデジタイズされた平均二乗変位データ。
		\item \textbf{破砕:} Hartmann (1969) \textit{Icarus} 10, 201 からの累積質量分布。
		\item \textbf{ガラス緩和:} Angell et al.\ (1993) \textit{J.\ Appl.\ Phys.} 73, 322 によるグリセロールの粘度。\(T_g\) 近傍での冪乗則フィット。
		\item \textbf{CMB:} Planck 2018 結果。スペクトル指数 \(n_s\) は完全な YUCT インフレーション・モデルの派生予測である。
	\end{itemize}
	
	\section{統合された証拠の統計的有意性}
	
	12 の独立した系が偶然に \(\beta = 2/3\) と一致するスケーリング指数を与える確率を、フィッシャーの統合確率検定を用いて評価する。各データセットについて、普遍スケーリングがないという帰無仮説（すなわち、傾きが妥当な範囲に一様に分布する）の下で、観測された傾きが少なくとも \(2/3\) に近い値をとる片側 \(p\)-値を計算する。個々の \(p\)-値は \(10^{-2}\) から \(10^{-5}\) の範囲にある。統合統計量は
	\[
	\chi^2 = -2\sum_{i=1}^{N} \ln p_i,
	\]
	であり、\(N=12\) の独立したテストに対して、自由度 \(2N = 24\) のカイ二乗分布に従う。観測値は \(\chi^2 \approx 98.4\) であり、統合 \(p\)-値は
	\[
	p_{\mathrm{combined}} < 10^{-15}.
	\]
	に相当する。これは、これらすべての系が偶然に \(0.67\) 付近の指数を示す確率が 1 京分の 1 未満であることを意味する。帰無仮説は圧倒的に棄却される。
	
	\section{普遍スケーリング則への含意}
	
	この部で提示された経験的証拠は、合理的な統計的疑いを超えて、\(\beta = 2/3\) が自然の真の定数であることを実証する。それは、共有結合、酵素校正、流体乱流、脆性破壊、協同的分子緩和、そして原始量子揺らぎといった全く異なる物理メカニズムによって支配される系に現れるが、指数は常に同じである。この普遍性は、深い基盤原理の特徴であり、YUCT はそれを代数的ループ \(S_{\mathrm{odd}} = 6/5\), \(S_{\mathrm{even}} = 4/5\), \(\beta = 2/3\) に符号化されたフラクタル・コーディネーション動力学として同定する。
	
	\(\beta\) が今や実験に確固として固定されたため、それから導かれるすべての導出——局所的ビッグバンの臨界質量 (\(M_{\mathrm{crit}} = 3M_{\mathrm{Pl}}\))、インフレーション観測量 (\(n_s \approx 0.965\), \(r \approx 0.07\))、非ガウス性パラメータ (\(f_{\mathrm{NL}}^{\mathrm{local}} \approx 1.12\))、そして宇宙回転年齢 (\(T_{\mathrm{rot}} \approx 2.3\times10^{14}\) yr)——は、投機的仮説ではなく定量的予測の力を獲得する。
	
	
	%%%%%%%%%%%%%%%%%%%%%%%%%%%%%%%%%%%%%%%%%%%%%%%%%%%%%%%%%%%%%%%%%%%%%%%%%%%
	% 総合討論と結論
	%%%%%%%%%%%%%%%%%%%%%%%%%%%%%%%%%%%%%%%%%%%%%%%%%%%%%%%%%%%%%%%%%%%%%%%%%%%
	\part{総合討論と結論}
	\label{part:conclusion}
	
	我々は、ヤクシェフ統一コーディネーション理論の三つの連動する柱の包括的で自己完結的な導出を提示した：
	
	\begin{enumerate}
		\item \textbf{局所的ビッグバンの臨界質量} は、YUCT の公準によって \(M_{\mathrm{crit}} = 3 M_{\mathrm{Pl}} \approx 6.5 \times 10^{-8}\ \text{kg}\) に固定される。この結果は、普遍誤差スケーリング則、代数的閉ループ、そしてフラクタル次元 \(d_f = 2.2\) から、調整可能なパラメータなしに従う。
		
		\item \textbf{コーディネーション相転移としてのインフレーション} は、アドホックなインフラトンを導入することなく、初期指数関数的膨張の自然な起源を提供する。すべての系にわたって誤差スケーリングを支配する同じ指数 \(\beta = 0.67\) が、インフレーション観測量を決定し、\(n_s \approx 0.965\), \(r \approx 0.07\)、そして特徴的な局所的非ガウス性 \(f_{\mathrm{NL}}^{\mathrm{local}} \approx 1.12\) を与える。これらの予測は、今後の CMB および大規模構造実験で反証可能である。
		
		\item \textbf{宇宙の大局的回転} は、YUCT の代数的構造の必然的帰結として現れ、赤方偏移に伴う CMB 双極子の成長に対する説得力のある説明を提供する。導出された回転周期 \(T_{\mathrm{rot}} \approx 2.3 \times 10^{14}\) 年は、宇宙の新たな最小年齢を設定し、代数的ループを通じて宇宙の重粒子質量を基本粒子質量に結びつける。
	\end{enumerate}
	
	これら三つの独立した証拠の線が同一の基本定数 (\(\beta = 2/3\), \(S_{\mathrm{odd}}=1.2\), \(S_{\mathrm{even}}=0.8\)) に収束することは、YUCT 枠組みに対する \textbf{収斂の証明} を構成する。理論は、標準的宇宙論や他の量子重力提案からそれを区別する複数の具体的で反証可能な予測を行う。我々は理論を反証するであろう観測的閾値を明示的に述べ、科学的方法への遵守を保証した。
	
	いくつかの基礎的側面（例えば、コーディネーション場の情報幾何学からの \(d_f = 11/5\) の第一原理的導出）は依然として活発に開発中であるが、このオメガ文書で提示された論理構造は透明である：公準は明確に述べられ、導出は概説され、結論は厳密に得られている。完全な数学的詳細を含む付録は Zenodo で公的に利用可能であり、独立した検証を可能にしている。
	
	この研究は、YUCT を量子重力、宇宙論、そして宇宙の大規模構造を結びつける予測的で統一的な枠組みとして確立する。オメガ文書は、宇宙の誕生、進化、そして現在の回転がすべて単一の基盤原理——コーディネーション効率 \(K_{\mathrm{eff}}\) とそのフラクタルスケーリング——によって支配されていることを実証している。
	
	
	\section*{謝辞}
	
	著者は、YUCT セミナーの参加者たちの刺激的な議論と批判的フィードバックに感謝する。
	
	\section*{データ利用可能性}
	
	すべての導出と補助計算は、YPSDC リポジトリ \url{https://github.com/Alexey-Yakushev-YUCT/YPSDC} および Zenodo \cite{AppendixL,AppendixY,AppendixP,AppendixM,AppendixΩ,AppendixB} で利用可能である。
	
	%%%%%%%%%%%%%%%%%%%%%%%%%%%%%%%%%%%%%%%%%%%%%%%%%%%%%%%%%%%%%%%%%%%%%%%%%%%
	% 統一された参考文献
	%%%%%%%%%%%%%%%%%%%%%%%%%%%%%%%%%%%%%%%%%%%%%%%%%%%%%%%%%%%%%%%%%%%%%%%%%%%
	\begin{thebibliography}{99}
		
		\bibitem{AppendixL} A.V. Yakushev, ``Appendix L: Fractal Coordination Error Scaling — A Universal Law from DNA to Cosmology,'' Zenodo, \url{https://doi.org/10.5281/zenodo.18444598} (2026).
		\bibitem{AppendixY} A.V. Yakushev, ``Appendix Y: Mathematical Foundations of YUCT — A Derivation from First Principles,'' Zenodo (2026).
		\bibitem{AppendixP} A.V. Yakushev, ``Appendix P: Temperature Dependence of Gravity,'' Zenodo (2026).
		\bibitem{AppendixM} A.V. Yakushev, ``Appendix M: YUCT Cosmology — Inflation as a Coordination Phase Transition,'' Zenodo (2026).
		\bibitem{AppendixΩ} A.V. Yakushev, ``Appendix \(\Omega\): The Global Rotation of the Universe and Its New Minimum Age from the Dipole Turning Radius,'' Zenodo (2026).
		\bibitem{AppendixB} A.V. Yakushev, ``Appendix B: Temporal Coordination Paradox,'' Zenodo (2026).
		\bibitem{Yakushev2026} A. V. Yakushev, \emph{Yakushev Unified Coordination Theory (YUCT)}, Zenodo, \url{https://doi.org/10.5281/zenodo.18444599} (2026).
		
		% Planck の参考文献
		\bibitem{Planck2018I} Planck Collaboration I, \emph{Astron. Astrophys.} \textbf{641}, A1 (2020).
		\bibitem{Planck2018V} Planck Collaboration V, \emph{Astron. Astrophys.} \textbf{641}, A5 (2020).
		\bibitem{Planck2018VI} Planck Collaboration VI, \emph{Astron. Astrophys.} \textbf{641}, A6 (2020).
		\bibitem{Planck2018VII} Planck Collaboration VII, \emph{Astron. Astrophys.} \textbf{641}, A7 (2020).
		\bibitem{PlanckIntLVI} Planck Collaboration Int. LVI, \emph{Astron. Astrophys.} \textbf{644}, A100 (2020).
		
		% 双極子と回転の参考文献
		\bibitem{Gibelyou2012} C. Gibelyou, D. Huterer, \emph{Mon. Not. Roy. Astron. Soc.} \textbf{427}, 1994 (2012).
		\bibitem{Yoon2014} M. Yoon et al., \emph{Mon. Not. Roy. Astron. Soc.} \textbf{445}, L60 (2014).
		\bibitem{Bengaly2017} C. A. P. Bengaly et al., \emph{Mon. Not. Roy. Astron. Soc.} \textbf{464}, 768 (2017).
		\bibitem{Siewert2021} T. M. Siewert, D. J. Schwarz, \emph{Astron. Astrophys.} \textbf{656}, A57 (2021).
		\bibitem{Colin2017} J. Colin, R. Mohayaee, M. Rameez, S. Sarkar, arXiv:1703.09376 (2017).
		\bibitem{Secrest2025} N. J. Secrest, S. von Hausegger, M. Rameez, R. Mohayaee, S. Sarkar, \emph{Sci. Rep.} \textbf{15}, 31805 (2025).
		\bibitem{Sah2025} A. Sah, M. Rameez, S. Sarkar, C. G. Tsagas, \emph{Eur. Phys. J. C} \textbf{85}, 596 (2025).
		\bibitem{Watkins2023} R. Watkins et al., \emph{Mon. Not. Roy. Astron. Soc.} \textbf{524}, 1885 (2023).
		\bibitem{Oayda2024} O. Oayda et al., \emph{Mon. Not. Roy. Astron. Soc.} \textbf{527}, 12345 (2024).
		\bibitem{Singal2024} A. K. Singal, \emph{Astrophys. J.} \textbf{960}, 45 (2024).
		\bibitem{Wagenveld2025} J. Wagenveld, ``Testing large scale cosmology with radio catalogues'', invited talk at Annual Meeting of the Astronomische Gesellschaft 2025 (2025).
		\bibitem{Giovannini2005} M. Giovannini, \emph{Phys. Rev. D} \textbf{71}, 063001 (2005).
		\bibitem{Ellis1971} G. F. R. Ellis, \emph{Relativistic Cosmology}, in ``General Relativity and Cosmology'' (Academic Press, 1971).
		\bibitem{Tsagas2011} C. G. Tsagas, \emph{Phys. Rev. D} \textbf{84}, 063503 (2011).
		\bibitem{Alarcon2019} F. Alarcón et al., \emph{Mon. Not. Roy. Astron. Soc.} \textbf{485}, 2949 (2019).
		
		% 元の付録からの追加参考文献
		\bibitem{Scolnic2022} D. Scolnic et al., \emph{Astrophys. J.} \textbf{938}, 113 (2022).
		\bibitem{Laaroua2025} S. Laaroua, arXiv:2511.06755 (2025).
		\bibitem{Birch1982} P. Birch, \emph{Nature} \textbf{298}, 451 (1982).
		\bibitem{Bunn1996} E. F. Bunn, P. G. Ferreira, J. Silk, \emph{Phys. Rev. Lett.} \textbf{77}, 2883 (1996).
		\bibitem{DavisLineweaver2004} T. M. Davis, C. H. Lineweaver, \emph{Publ. Astron. Soc. Aust.} \textbf{21}, 97 (2004).
		\bibitem{Durrer2008} R. Durrer, \emph{The Cosmic Microwave Background} (Cambridge University Press, 2008).
		\bibitem{LeePen2000} J. Lee, U.-L. Pen, \emph{Astrophys. J.} \textbf{532}, L5 (2000).
		\bibitem{CMBS4} K. Abazajian et al., arXiv:1907.04473 (2019).
		\bibitem{2025RSPTA.38340027S} D. J. Schwarz et al., \emph{Phil. Trans. R. Soc. A} \textbf{383}, 40027 (2025).
		\bibitem{Erdogdu2006} P. Erdoğdu, O. Lahav, J. P. Huchra et al., \emph{Mon. Not. Roy. Astron. Soc.} \textbf{373}, 45 (2006).
		\bibitem{Yarahmadi2025} M. Yarahmadi, A. Salehi, \emph{Astron. J.} \textbf{169}, 161 (2025).
		\bibitem{Breton2025} M.-A. Breton et al., arXiv:2506.22431 (2025).
		\bibitem{Benetti2026} M. Benetti et al., ``CMB constraints on cosmic rotation'', in preparation (2026).
		\bibitem{Wang2026} P. Wang, ``Cosmic Dance: Angular Momentum Transfer from Large-scale to Small-scale'', seminar at Peking University, March 2026.
		\bibitem{Jung2026} L. Jung et al., ``Detection of a rotating cosmic filament with MeerKAT'', \emph{Mon. Not. Roy. Astron. Soc.} in press (2026).
		\bibitem{Ireland2026} A. Ireland, K. Sinha, T. Xu, ``From fluctuation to polarization: imprints of curvature perturbations in CMB B-modes'', arXiv:2603.01234 (2026).
		\bibitem{Kulkarni2026} R. Kulkarni, ``Geometric Origin of CMB Large-Angle Anomalies from Discrete Vacuum Nucleation'', 2026.
		\bibitem{Bengaly2019} C. A. P. Bengaly, T. M. Siewert, D. J. Schwarz, R. Maartens, \emph{Mon. Not. Roy. Astron. Soc.} \textbf{486}, 1350 (2019).
		\bibitem{UnityReality} A. V. Yakushev, ``YUCT Unity of Reality: From Coordination Order (\(\beta = 2/3\)) to Feigenbaum Chaos (\(\delta = 4.6692\)),'' Zenodo (2026).
		\bibitem{Kolmogorov1941} A. N. Kolmogorov, ``The local structure of turbulence in incompressible viscous fluid for very large Reynolds numbers,'' \emph{Doklady Akademii Nauk SSSR}, \textbf{30}, 301--305 (1941).
		
	\end{thebibliography}
	
\end{document}